% CP-VI-0047
% target_chapter: VI/22 — Open and Closed Subschemes
% source_unit_id: IMP-DL-0A9B1F0A75-P004
% source: Downloads/theory-of-algebraic-geometry-9.tex L1416-L5921
% solution_provenance: SOURCE_EXPLICIT_SOLUTION

\subsection*{Problem statement}

Let
\[
i:Z\hookrightarrow X
\]
be a closed immersion of schemes. Recall that this means that \(i\) is a homeomorphism of \(Z\) onto a closed subset of \(X\), and that the morphism of sheaves
\[
i^\#:\mathcal{O}_X\longrightarrow i_*\mathcal{O}_Z
\]
is surjective.

Define
\[
\mathcal{I}_{Z/X}=\ker i^\#.
\]

Show the following.

\begin{enumerate}[label=(\alph*)]
	\item \(\mathcal{I}_{Z/X}\) is a sheaf of ideals of \(\mathcal{O}_X\); that is,
	\[
	\mathcal{I}_{Z/X}(U)
	\]
	is an ideal of \(\mathcal{O}_X(U)\) for every open subset \(U\subseteq X\).
	
	\item \(\mathcal{I}_{Z/X}\) is a quasi-coherent sheaf of \(\mathcal{O}_X\)-modules, and it is coherent if \(X\) is Noetherian.
	
	\item There is a one-to-one correspondence between quasi-coherent sheaves of ideals on \(X\) and closed subschemes of \(X\).
\end{enumerate}

% ------------------------------------------------------------
\subsection{\texorpdfstring{\(\mathcal{I}_{Z/X}\)}{I Z/X} is a sheaf of ideals}

\textbf{Solution.}

By definition,
\[
\mathcal{I}_{Z/X}
=
\ker\!\left(
\mathcal{O}_X\longrightarrow i_*\mathcal{O}_Z
\right).
\]

For every open subset
\[
U\subseteq X,
\]
we have
\[
(i_*\mathcal{O}_Z)(U)
=
\mathcal{O}_Z(i^{-1}(U)).
\]

Since \(i\) identifies \(Z\) homeomorphically with a closed subset of \(X\), we may write
\[
i^{-1}(U)=U\cap Z.
\]

Therefore
\[
\mathcal{I}_{Z/X}(U)
=
\ker\!\left(
\mathcal{O}_X(U)
\longrightarrow
\mathcal{O}_Z(U\cap Z)
\right).
\]

The map
\[
\mathcal{O}_X(U)\longrightarrow \mathcal{O}_Z(U\cap Z)
\]
is a ring homomorphism. The kernel of a ring homomorphism is an ideal. Hence
\[
\mathcal{I}_{Z/X}(U)
\]
is an ideal in
\[
\mathcal{O}_X(U).
\]

Now let
\[
V\subseteq U
\]
be open subsets of \(X\). The restriction maps fit into a commutative diagram
\[
\begin{array}{ccc}
	\mathcal{O}_X(U) & \longrightarrow & \mathcal{O}_Z(U\cap Z) \\[4pt]
	\downarrow & & \downarrow \\[4pt]
	\mathcal{O}_X(V) & \longrightarrow & \mathcal{O}_Z(V\cap Z).
\end{array}
\]

Therefore restriction maps send elements of
\[
\mathcal{I}_{Z/X}(U)
\]
to elements of
\[
\mathcal{I}_{Z/X}(V).
\]

Thus \(\mathcal{I}_{Z/X}\) is a subsheaf of rings of \(\mathcal{O}_X\), and on every open set its sections form an ideal.

Hence
\[
\boxed{
	\mathcal{I}_{Z/X}
	\text{ is a sheaf of ideals of }\mathcal{O}_X.
}
\]

% ------------------------------------------------------------
\subsection{\texorpdfstring{\(\mathcal{I}_{Z/X}\)}{I Z/X} is quasi-coherent and, in the Noetherian case, coherent}

\textbf{Solution.}

Take an affine open subset
\[
U=\operatorname{Spec} A\subseteq X.
\]

Because
\[
i:Z\hookrightarrow X
\]
is a closed immersion, the intersection
\[
Z\cap U
\]
is a closed subscheme of \(U\). Closed subschemes of an affine scheme correspond to ideals. Therefore there exists an ideal
\[
I\subseteq A
\]
such that
\[
Z\cap U\cong \operatorname{Spec}(A/I).
\]

On \(U\), the morphism
\[
\mathcal{O}_X\longrightarrow i_*\mathcal{O}_Z
\]
is induced by the quotient map
\[
A\longrightarrow A/I.
\]

Its kernel is the ideal \(I\). Passing to associated sheaves, we obtain
\[
\mathcal{I}_{Z/X}|_U
\cong
\widetilde{I}.
\]

Thus, on every affine open subset \(U=\operatorname{Spec} A\subseteq X\), the ideal sheaf \(\mathcal{I}_{Z/X}\) is associated to an \(A\)-module. Hence it is quasi-coherent:
\[
\boxed{
	\mathcal{I}_{Z/X}
	\text{ is a quasi-coherent sheaf of }\mathcal{O}_X\text{-modules.}
}
\]

Now assume that \(X\) is Noetherian.

Then for every affine open subset
\[
U=\operatorname{Spec} A\subseteq X,
\]
the ring \(A\) is Noetherian. In particular, every ideal
\[
I\subseteq A
\]
is finitely generated.

Therefore
\[
\widetilde{I}
\]
is a coherent sheaf on \(U\).

Since coherence is local, it follows that
\[
\boxed{
	X\text{ Noetherian}
	\quad\Longrightarrow\quad
	\mathcal{I}_{Z/X}
	\text{ is coherent.}
}
\]

% ------------------------------------------------------------
\subsection{Correspondence between closed subschemes and quasi-coherent ideal sheaves}

\textbf{Solution.}

We shall prove the bijection
\[
\boxed{
	\left\{
	\begin{array}{c}
		\text{closed subschemes}\\
		Z\hookrightarrow X
	\end{array}
	\right\}
	\longleftrightarrow
	\left\{
	\begin{array}{c}
		\text{quasi-coherent sheaves of ideals}\\
		\mathcal{I}\subseteq \mathcal{O}_X
	\end{array}
	\right\}.
}
\]

\subsubsection{From a closed subscheme to a quasi-coherent ideal sheaf}

Let
\[
i:Z\hookrightarrow X
\]
be a closed immersion.

Define
\[
\mathcal{I}_{Z/X}
=
\ker\!\left(
\mathcal{O}_X\longrightarrow i_*\mathcal{O}_Z
\right).
\]

By the previous parts:

\begin{itemize}
	\item \(\mathcal{I}_{Z/X}\) is a sheaf of ideals of \(\mathcal{O}_X\);
	\item \(\mathcal{I}_{Z/X}\) is quasi-coherent.
\end{itemize}

Therefore every closed subscheme determines a quasi-coherent ideal sheaf.

\subsubsection{From a quasi-coherent ideal sheaf to a closed subscheme}

Conversely, let
\[
\mathcal{I}\subseteq \mathcal{O}_X
\]
be a quasi-coherent sheaf of ideals.

Take an affine open subset
\[
U=\operatorname{Spec} A\subseteq X.
\]

Since \(\mathcal{I}\) is quasi-coherent, there exists an ideal
\[
I_U\subseteq A
\]
such that
\[
\mathcal{I}|_U\cong \widetilde{I_U}.
\]

Define a closed subscheme of \(U\) by
\[
Z_U=\operatorname{Spec}(A/I_U).
\]

Its underlying closed subset is
\[
V(I_U)\subseteq \operatorname{Spec} A.
\]

Now suppose
\[
V=\operatorname{Spec} B\subseteq U
\]
is another affine open subset. Since the ideal sheaf restricts compatibly, the ideal defining \(Z_V\) is the localization of \(I_U\) corresponding to \(V\). Consequently,
\[
Z_V
\]
is precisely the restriction of
\[
Z_U
\]
to \(V\).

Thus the affine closed subschemes
\[
Z_U=\operatorname{Spec}(A/I_U)
\]
agree on overlaps and glue together to form a scheme \(Z\) together with a closed immersion
\[
i:Z\hookrightarrow X.
\]

Locally, the structure sheaf of \(Z\) is
\[
\mathcal{O}_Z|_{Z\cap U}
\cong
\widetilde{A/I_U}.
\]

Equivalently, the quotient sheaf
\[
\mathcal{O}_X/\mathcal{I}
\]
is the sheaf of rings defining the closed subscheme \(Z\).

The morphism
\[
\mathcal{O}_X\longrightarrow i_*\mathcal{O}_Z
\]
has kernel exactly \(\mathcal{I}\). Hence
\[
\mathcal{I}_{Z/X}=\mathcal{I}.
\]

\subsubsection{The constructions are inverse}

Starting with a closed subscheme
\[
Z\hookrightarrow X,
\]
forming its ideal sheaf
\[
\mathcal{I}_{Z/X},
\]
and then reconstructing the closed subscheme locally by quotient rings gives back the original closed subscheme \(Z\).

Conversely, starting with a quasi-coherent ideal sheaf
\[
\mathcal{I}\subseteq \mathcal{O}_X,
\]
forming the local closed subschemes
\[
\operatorname{Spec}(A/I_U),
\]
and then taking the ideal sheaf of the resulting closed immersion gives back \(\mathcal{I}\).

Therefore the constructions are mutually inverse, and we obtain the required bijection:
\[
\boxed{
	\left\{
	\text{closed subschemes of }X
	\right\}
	\longleftrightarrow
	\left\{
	\text{quasi-coherent ideal sheaves }
	\mathcal{I}\subseteq \mathcal{O}_X
	\right\}.
}
\]

\subsection{The affine special case}

When
\[
X=\operatorname{Spec} A
\]
is affine, the correspondence becomes especially concrete.

Every ideal
\[
I\subseteq A
\]
defines a closed subscheme
\[
\operatorname{Spec}(A/I)\hookrightarrow \operatorname{Spec} A.
\]

Conversely, every closed subscheme of \(\operatorname{Spec} A\) arises uniquely in this way.

Thus
\[
\boxed{
	\left\{
	\text{closed subschemes of }\operatorname{Spec} A
	\right\}
	\longleftrightarrow
	\left\{
	\text{ideals }I\subseteq A
	\right\},
	\qquad
	I\longmapsto \operatorname{Spec}(A/I).
}
\]

This is the scheme-theoretic refinement of the ordinary correspondence between ideals of \(A\) and closed subsets of \(\operatorname{Spec} A\): different ideals may define the same underlying closed set, but they define different closed subscheme structures.

For example,
\[
(x)\subseteq k[x]
\qquad\text{and}\qquad
(x^2)\subseteq k[x]
\]
both define the same underlying closed point
\[
\{(x)\}\subseteq \mathbf{A}^1_k,
\]
but the closed subschemes
\[
\operatorname{Spec} k[x]/(x)
\qquad\text{and}\qquad
\operatorname{Spec} k[x]/(x^2)
\]
are different. The second is a nonreduced thickening of the first.

% ============================================================

\subsection*{Theory Behind Problems 4--7:
	Divisor Classes, Projective Morphisms, Coherent Sheaves, and Closed Subschemes}

The problems above bring together four fundamental themes of scheme theory:

\begin{enumerate}[label=\arabic*.]
	\item divisor class groups measure the failure of codimension-one subvarieties to be principal;
	\item morphisms from projective space are strongly constrained by positivity and dimension;
	\item quasi-coherent and coherent sheaves translate geometric constructions into module-theoretic constructions on affine charts;
	\item closed subschemes are encoded exactly by quasi-coherent ideal sheaves.
\end{enumerate}

Although these topics appear different, they are connected by one central principle:

\begin{center}
	\emph{Global geometric questions on schemes are often solved by passing to affine open subsets, performing algebra there, and then gluing the conclusions globally.}
\end{center}

% ------------------------------------------------------------
\subsection{Divisors and divisor class groups on normal schemes}

\subsubsection{Codimension-one geometry}

Let \(X\) be an integral Noetherian scheme. A closed irreducible subset
\[
D\subseteq X
\]
of codimension one is called a \emph{prime divisor}.

Prime divisors are the scheme-theoretic analogues of irreducible hypersurfaces. On an affine variety, they are usually cut out locally by a single irreducible equation. However, globally a codimension-one subvariety need not be defined by one global function.

This difference is precisely what the divisor class group detects.

A \emph{Weil divisor} on \(X\) is a finite formal integral linear combination of prime divisors:
\[
D=\sum_{Y\subseteq X,\ \operatorname{codim}(Y,X)=1} n_Y[Y],
\qquad
n_Y\in \mathbb{Z},
\]
where all but finitely many \(n_Y\) are zero.

The abelian group of Weil divisors is denoted
\[
\operatorname{Div}(X).
\]

Thus a divisor records codimension-one vanishing and pole data. For example, if
\[
D=2D_1-D_2,
\]
then \(D\) should be interpreted as having multiplicity \(2\) along \(D_1\) and multiplicity \(-1\), that is, a pole, along \(D_2\).

\subsubsection{Why normality is important}

Suppose that \(X\) is normal and integral. Then every local ring
\[
\mathcal{O}_{X,\eta_D}
\]
at the generic point \(\eta_D\) of a prime divisor \(D\) is a discrete valuation ring.

Consequently, every nonzero rational function
\[
f\in K(X)^\times
\]
has a well-defined order of vanishing along \(D\):
\[
\operatorname{ord}_D(f)\in \mathbb{Z}.
\]

The integer \(\operatorname{ord}_D(f)\) has the following meaning:

\[
\operatorname{ord}_D(f)>0
\quad\Longleftrightarrow\quad
f\text{ vanishes along }D,
\]
\[
\operatorname{ord}_D(f)<0
\quad\Longleftrightarrow\quad
f\text{ has a pole along }D,
\]
and
\[
\operatorname{ord}_D(f)=0
\quad\Longleftrightarrow\quad
f\text{ is a unit at the generic point of }D.
\]

The \emph{principal divisor} associated to \(f\) is
\[
\operatorname{div}(f)
=
\sum_D \operatorname{ord}_D(f)[D].
\]

Only finitely many coefficients are nonzero, so this is a Weil divisor.

\begin{definition}
	Let \(X\) be a normal integral Noetherian scheme. The \emph{divisor class group} of \(X\) is
	\[
	\operatorname{Cl}(X)
	=
	\frac{\operatorname{Div}(X)}
	{\{\operatorname{div}(f):f\in K(X)^\times\}}.
	\]
\end{definition}

Thus two divisors represent the same class precisely when their difference is the divisor of a rational function.

A divisor class is zero if and only if the divisor is principal.

\subsubsection{Geometric meaning of the class group}

The group \(\operatorname{Cl}(X)\) measures the obstruction to describing codimension-one geometry by global rational functions.

If
\[
\operatorname{Cl}(X)=0,
\]
then every Weil divisor is principal. In an affine normal setting this is closely related to unique factorization.

For a normal domain \(A\),
\[
\operatorname{Cl}(\operatorname{Spec} A)=0
\]
if \(A\) is a unique factorization domain. Indeed, height-one prime ideals correspond to irreducible codimension-one subvarieties, and unique factorization forces each such divisor to arise from one equation.

Thus:

\[
\boxed{
	A\text{ a normal UFD}
	\quad\Longrightarrow\quad
	\operatorname{Cl}(\operatorname{Spec} A)=0.
}
\]

This explains why polynomial rings and Laurent polynomial rings are so useful in class-group calculations:
\[
k[x_1,\dots,x_n],
\qquad
k[x_1^{\pm 1},x_2,\dots,x_n]
\]
are UFDs, so their spectra have trivial class group.

\subsubsection{Removing divisors and the localization sequence}

A powerful method for computing divisor class groups is to remove a closed subset until the remaining open subscheme becomes factorial.

Let \(X\) be normal and integral, and let
\[
U=X\setminus \bigcup_{i=1}^{r}D_i,
\]
where the \(D_i\) are precisely the irreducible codimension-one components of \(X\setminus U\).

Then restriction of divisors gives an exact sequence
\[
\bigoplus_{i=1}^{r}\mathbb{Z}[D_i]
\longrightarrow
\operatorname{Cl}(X)
\longrightarrow
\operatorname{Cl}(U)
\longrightarrow 0.
\]

The meaning of this sequence is simple:

\begin{itemize}
	\item divisors on \(X\) restrict to divisors on \(U\);
	\item divisors supported entirely in \(X\setminus U\) disappear after restriction;
	\item the only codimension-one divisors that disappear are the removed divisors \(D_i\).
\end{itemize}

If \(U\) is factorial, then
\[
\operatorname{Cl}(U)=0.
\]

Hence the removed divisors generate \(\operatorname{Cl}(X)\):
\[
\operatorname{Cl}(X)
=
\langle [D_1],\dots,[D_r]\rangle.
\]

The remaining question is to find all relations among those generators. These relations arise from principal divisors of rational functions whose zero and pole loci lie in the removed boundary.

\subsubsection{Units on the factorial open subset}

Suppose
\[
\operatorname{div}(f)
\]
is supported entirely in \(X\setminus U\). Then on \(U\), the rational function \(f\) has no zeros and no poles. Therefore
\[
f|_U\in \Gamma(U,\mathcal{O}_U)^\times.
\]

This converts a geometric relation in \(\operatorname{Cl}(X)\) into a computation of units in a ring.

For example:
\[
k[s,t]^\times=k^\times,
\]
whereas
\[
k[x,x^{-1},z]^\times=k^\times x^{\mathbb{Z}}.
\]

These two unit computations lead to two different types of behavior:

\begin{itemize}
	\item if the only units are constants, then there are no nontrivial divisor relations supported on the boundary;
	\item if a Laurent monomial such as \(x^m\) is also a unit, then powers of \(x\) produce relations among boundary divisors.
\end{itemize}

This is the basic mechanism behind the three computations in Problem 4.

% ------------------------------------------------------------
\subsection{The class group of a product of projective lines}

The surface
\[
\mathbf{P}^1_k\times \mathbf{P}^1_k
\]
is the simplest example of a projective surface carrying two independent families of curves.

The two projections
\[
\pi_1:\mathbf{P}^1_k\times\mathbf{P}^1_k\longrightarrow\mathbf{P}^1_k,
\qquad
\pi_2:\mathbf{P}^1_k\times\mathbf{P}^1_k\longrightarrow\mathbf{P}^1_k
\]
produce two natural divisor classes:
\[
F_1=\{\text{point}\}\times \mathbf{P}^1_k,
\qquad
F_2=\mathbf{P}^1_k\times\{\text{point}\}.
\]

These are called the two \emph{rulings} of the surface.

Removing one representative from each ruling gives
\[
(\mathbf{P}^1_k\times\mathbf{P}^1_k)
\setminus
\left(
(\{\infty\}\times\mathbf{P}^1_k)
\cup
(\mathbf{P}^1_k\times\{\infty\})
\right)
\cong
\mathbf{A}^2_k.
\]

Since
\[
\operatorname{Cl}(\mathbf{A}^2_k)=0,
\]
the two boundary divisors generate the class group.

What distinguishes this example from the affine cones is that
\[
\Gamma(\mathbf{A}^2_k,\mathcal{O}_{\mathbf{A}^2_k})^\times
=
k^\times.
\]

Thus no nonconstant rational function can have its entire divisor supported only on these two boundary curves. Hence the two ruling classes remain independent:
\[
\operatorname{Cl}(\mathbf{P}^1_k\times\mathbf{P}^1_k)\cong \mathbb{Z}^2.
\]

There is also a line-bundle interpretation. Since the surface is smooth, every Weil divisor is Cartier, and therefore
\[
\operatorname{Cl}(\mathbf{P}^1_k\times\mathbf{P}^1_k)
\cong
\operatorname{Pic}(\mathbf{P}^1_k\times\mathbf{P}^1_k).
\]

The corresponding line bundles are
\[
\mathcal{O}_{\mathbf{P}^1\times\mathbf{P}^1}(a,b)
=
\pi_1^*\mathcal{O}_{\mathbf{P}^1}(a)\otimes
\pi_2^*\mathcal{O}_{\mathbf{P}^1}(b),
\]
and every line bundle is of this form. Hence
\[
\operatorname{Pic}(\mathbf{P}^1_k\times\mathbf{P}^1_k)
\cong
\mathbb{Z}\oplus \mathbb{Z}.
\]

The pair \((a,b)\) records the degree of the divisor in the two ruling directions.

% ------------------------------------------------------------
\subsection{Quadric cones and nontrivial divisor classes}

\subsubsection{The surface cone \(\boldsymbol{xy=z^2}\)}

Consider
\[
X=\operatorname{Spec} k[x,y,z]/(xy-z^2).
\]

This is a two-dimensional affine cone. Away from the origin it is smooth, but the vertex creates nontrivial divisor-class behavior.

If we invert \(x\), then
\[
y=\frac{z^2}{x},
\]
so
\[
\mathcal{O}_X(D(x))
\cong
k[x,x^{-1},z].
\]

This ring is factorial. Therefore all divisor-class information is concentrated in the closed set
\[
V(x).
\]

Scheme-theoretically,
\[
A/(x)\cong k[y,z]/(z^2).
\]

The underlying irreducible codimension-one subset is
\[
D=V(x,z).
\]

The important point is that \(x\) does not vanish simply along \(D\). At the generic point of \(D\), the element \(y\) is a unit and
\[
x=\frac{z^2}{y}.
\]

Thus \(x\) vanishes to order \(2\) along \(D\):
\[
\operatorname{div}(x)=2D.
\]

Hence the class \([D]\) is not necessarily zero, but twice the class is zero:
\[
2[D]=0.
\]

This is the geometric origin of the torsion class
\[
\operatorname{Cl}(X)\cong \mathbb{Z}/2\mathbb{Z}.
\]

The cone has a divisor which behaves like ``half'' of a principal divisor: the function \(x\) defines twice the ruling divisor.

\subsubsection{The three-dimensional cone \(\boldsymbol{xy=zw}\)}

Now consider
\[
X=\operatorname{Spec} k[x,y,z,w]/(xy-zw).
\]

Again inverting \(x\) makes the ring factorial:
\[
A_x\cong k[x,x^{-1},z,w].
\]

However, the boundary \(V(x)\) now has two irreducible codimension-one components:
\[
D=V(x,z),
\qquad
E=V(x,w).
\]

At the generic point of \(D\),
\[
x=\frac{w}{y}z,
\]
so \(x\) vanishes to order \(1\) along \(D\). Similarly, at the generic point of \(E\),
\[
x=\frac{z}{y}w,
\]
so \(x\) vanishes to order \(1\) along \(E\).

Thus
\[
\operatorname{div}(x)=D+E.
\]

This gives the relation
\[
[D]+[E]=0.
\]

Unlike the surface cone \(xy=z^2\), there is no multiplicity \(2\) along a single divisor. Instead, one principal divisor splits into two different prime divisors. Therefore one divisor class remains freely generated:
\[
\operatorname{Cl}(X)\cong \mathbb{Z}.
\]

Conceptually, the comparison is:

\[
\begin{array}{c|c|c}
	\textbf{Scheme} & \textbf{Boundary divisors} & \textbf{Principal relation} \\
	\hline
	\operatorname{Spec} k[x,y,z]/(xy-z^2)
	&
	D=V(x,z)
	&
	\operatorname{div}(x)=2D
	\\[4pt]
	\operatorname{Spec} k[x,y,z,w]/(xy-zw)
	&
	D=V(x,z),\ E=V(x,w)
	&
	\operatorname{div}(x)=D+E
\end{array}
\]

The first relation produces torsion:
\[
\mathbb{Z}/(2)\cong \mathbb{Z}/2\mathbb{Z}.
\]

The second relation removes one of two free generators:
\[
\mathbb{Z}^2/\mathbb{Z}(1,1)\cong \mathbb{Z}.
\]

% ------------------------------------------------------------
\subsection{Morphisms from projective space and positivity}

\subsubsection{Morphisms to projective space as homogeneous polynomials}

A morphism
\[
\varphi:\mathbf{P}^n\longrightarrow\mathbf{P}^m
\]
is controlled by a line bundle and a collection of global sections.

Indeed, let
\[
\mathcal{O}_{\mathbf{P}^m}(1)
\]
be the hyperplane line bundle on \(\mathbf{P}^m\). Its pullback is a line bundle on \(\mathbf{P}^n\):
\[
\varphi^*\mathcal{O}_{\mathbf{P}^m}(1).
\]

Since
\[
\operatorname{Pic}(\mathbf{P}^n)\cong\mathbb{Z},
\]
every line bundle on \(\mathbf{P}^n\) is of the form
\[
\mathcal{O}_{\mathbf{P}^n}(d)
\]
for a unique integer \(d\). Since \(\varphi^*\mathcal{O}_{\mathbf{P}^m}(1)\) is generated by global sections, one necessarily has
\[
d\geq 0.
\]

Thus
\[
\varphi^*\mathcal{O}_{\mathbf{P}^m}(1)\cong \mathcal{O}_{\mathbf{P}^n}(d)
\]
for some \(d\geq 0\).

The morphism \(\varphi\) is then given by \(m+1\) global sections of
\[
\mathcal{O}_{\mathbf{P}^n}(d),
\]
that is, by homogeneous polynomials
\[
F_0,\dots,F_m
\]
of the same degree \(d\), with no common zero:
\[
\varphi([X_0:\cdots:X_n])
=
[F_0(X):\cdots:F_m(X)].
\]

The condition that the \(F_i\) have no common zero is essential. If they vanished simultaneously at a point, then the projective coordinate expression would be undefined there.

\subsubsection{Constant maps and degree zero}

If
\[
d=0,
\]
then each \(F_i\) is a constant. Therefore the image consists of a single projective point:
\[
\varphi(\mathbf{P}^n)=\{[c_0:\cdots:c_m]\}.
\]

Thus:
\[
\boxed{
	\varphi^*\mathcal{O}_{\mathbf{P}^m}(1)\cong \mathcal{O}_{\mathbf{P}^n}
	\quad\Longrightarrow\quad
	\varphi\text{ is constant.}
}
\]

Conversely, a constant map clearly pulls back \(\mathcal{O}_{\mathbf{P}^m}(1)\) to the trivial line bundle.

Hence nonconstant morphisms correspond to the case
\[
d>0.
\]

\subsubsection{Why there are no nonconstant morphisms \(\boldsymbol{\mathbf{P}^n\to\mathbf{P}^r}\) for \(\boldsymbol{r<n}\)}

Suppose
\[
\psi:\mathbf{P}^n\longrightarrow\mathbf{P}^r
\]
is nonconstant, where
\[
r<n.
\]

Then \(\psi\) is given by \(r+1\) homogeneous polynomials
\[
F_0,\dots,F_r
\]
of a positive common degree.

For these polynomials to define a morphism, they must have no common zero in \(\mathbf{P}^n\).

However, \(r+1\leq n\). Geometrically, fewer than \(n+1\) positive-degree hypersurfaces in \(\mathbf{P}^n\) cannot cut out the empty set. Their intersection must contain at least one projective point.

Algebraically, the ideal
\[
(F_0,\dots,F_r)
\subseteq
k[X_0,\dots,X_n]
\]
has height at most
\[
r+1\leq n,
\]
whereas the irrelevant ideal
\[
(X_0,\dots,X_n)
\]
has height
\[
n+1.
\]

Thus the radical of \((F_0,\dots,F_r)\) cannot contain the irrelevant ideal. Therefore the \(F_i\) possess a common projective zero, contradicting the assumption that they define a morphism everywhere.

Hence:
\[
\boxed{
	r<n
	\quad\Longrightarrow\quad
	\text{every morphism }\mathbf{P}^n\longrightarrow\mathbf{P}^r
	\text{ is constant.}
}
\]

This is a particularly strong rigidity property of projective space.

\subsubsection{Projection from a linear subspace}

Suppose now that
\[
\varphi:\mathbf{P}^n\longrightarrow\mathbf{P}^m
\]
has image
\[
Y=\varphi(\mathbf{P}^n)
\]
of dimension
\[
r.
\]

If
\[
0<r<n,
\]
one can choose a sufficiently general linear subspace
\[
\Lambda\subseteq\mathbf{P}^m
\]
of dimension
\[
m-r-1
\]
which does not meet \(Y\).

Projection away from \(\Lambda\) gives a morphism
\[
\pi_{\Lambda}:Y\longrightarrow\mathbf{P}^r.
\]

For general \(\Lambda\), this projection is finite onto its image. In particular, it is nonconstant whenever
\[
r>0.
\]

The composite
\[
\mathbf{P}^n
\xrightarrow{\varphi}
Y
\xrightarrow{\pi_{\Lambda}}
\mathbf{P}^r
\]
would therefore be a nonconstant morphism from \(\mathbf{P}^n\) to \(\mathbf{P}^r\) with
\[
r<n.
\]

But such a morphism cannot exist. Consequently, the dimension of the image cannot satisfy
\[
0<\dim \varphi(\mathbf{P}^n)<n.
\]

Therefore the image is either zero-dimensional or full-dimensional:
\[
\boxed{
	\dim\varphi(\mathbf{P}^n)=0
	\quad\text{or}\quad
	\dim\varphi(\mathbf{P}^n)=n.
}
\]

Since an irreducible zero-dimensional projective variety is a point, this becomes:
\[
\boxed{
	\varphi(\mathbf{P}^n)\text{ is a point, or }
	\dim\varphi(\mathbf{P}^n)=n.
}
\]

In the nonconstant case,
\[
\varphi(\mathbf{P}^n)\subseteq\mathbf{P}^m
\]
has dimension \(n\), so necessarily
\[
m\geq n.
\]

% ------------------------------------------------------------
\subsection{Quasi-coherent sheaves as geometric modules}

\subsubsection{From modules to sheaves}

Let
\[
X=\operatorname{Spec} A
\]
be an affine scheme and let \(M\) be an \(A\)-module. Associated to \(M\) is a sheaf of \(\mathcal{O}_X\)-modules
\[
\widetilde{M}.
\]

Its defining property on distinguished open subsets is
\[
\widetilde{M}(D(f))=M_f.
\]

Thus localization of modules becomes restriction of sheaves.

For example:
\[
\widetilde{A}=\mathcal{O}_X,
\]
and for an ideal \(I\subseteq A\),
\[
\widetilde{I}\subseteq \mathcal{O}_X
\]
is the corresponding ideal sheaf.

\begin{definition}
	A sheaf \(\mathcal{F}\) of \(\mathcal{O}_X\)-modules is \emph{quasi-coherent} if every point has an affine open neighbourhood
	\[
	U=\operatorname{Spec} A
	\]
	such that
	\[
	\mathcal{F}|_U\cong\widetilde{M}
	\]
	for some \(A\)-module \(M\).
\end{definition}

Quasi-coherent sheaves are therefore the sheaves that locally come from modules.

This makes them the natural class of sheaves in algebraic geometry. They retain enough algebraic structure to be manipulated locally, while being flexible enough to glue across different affine charts.

\subsubsection{The affine equivalence}

For an affine scheme
\[
X=\operatorname{Spec} A,
\]
there is an equivalence of categories
\[
\left\{
A\text{-modules}
\right\}
\longleftrightarrow
\left\{
\text{quasi-coherent }\mathcal{O}_X\text{-modules}
\right\},
\]
given by
\[
M\longmapsto \widetilde M
\]
and
\[
\mathcal{F}\longmapsto \Gamma(X,\mathcal{F}).
\]

This means that on affine schemes, questions about quasi-coherent sheaves are exactly questions about modules.

In particular, if
\[
\phi:M\longrightarrow N
\]
is a module homomorphism, then
\[
\widetilde{\phi}:\widetilde{M}\longrightarrow\widetilde{N}
\]
is the corresponding sheaf morphism.

Because localization is exact, one has:
\[
\ker(\widetilde{\phi})
\cong
\widetilde{\ker\phi},
\]
\[
\operatorname{im}(\widetilde{\phi})
\cong
\widetilde{\operatorname{im}\phi},
\]
and
\[
\operatorname{coker}(\widetilde{\phi})
\cong
\widetilde{\operatorname{coker}\phi}.
\]

This is the local algebraic reason that kernels, images, and cokernels of morphisms of quasi-coherent sheaves remain quasi-coherent.

\subsubsection{Abelian-category behavior}

The category of all sheaves of \(\mathcal{O}_X\)-modules is an abelian category: kernels, cokernels, and images always exist as sheaves.

The nontrivial statement is that when one begins with quasi-coherent sheaves, these constructions do not leave the quasi-coherent world:
\[
\mathcal{F},\mathcal{G}\text{ quasi-coherent}
\quad\Longrightarrow\quad
\ker f,\ \operatorname{im} f,\ \operatorname{coker} f\text{ quasi-coherent}.
\]

Therefore quasi-coherent sheaves form an abelian subcategory of the category of all \(\mathcal{O}_X\)-modules.

This is fundamental because exact sequences such as
\[
0\longrightarrow\mathcal{F}'
\longrightarrow\mathcal{F}
\longrightarrow\mathcal{F}''
\longrightarrow 0
\]
can be studied locally as exact sequences of modules:
\[
0\longrightarrow M'
\longrightarrow M
\longrightarrow M''
\longrightarrow 0.
\]

% ------------------------------------------------------------
\subsection{Coherent sheaves and finiteness}

\subsubsection{Finite generation and finite presentation}

Quasi-coherent sheaves may be extremely large. For instance, on
\[
\operatorname{Spec} k
\]
every \(k\)-vector space defines a quasi-coherent sheaf, including infinite-dimensional ones.

Coherent sheaves are the finite objects among quasi-coherent sheaves.

On a locally Noetherian scheme, a quasi-coherent sheaf \(\mathcal{F}\) is coherent precisely when, on every affine open subset
\[
U=\operatorname{Spec} A,
\]
one has
\[
\mathcal{F}|_U\cong\widetilde{M}
\]
for a finitely generated \(A\)-module \(M\).

Since \(A\) is Noetherian, finitely generated modules have finitely generated submodules. Therefore kernels and images remain finitely generated.

This is exactly the point at which the Noetherian hypothesis becomes important.

\subsubsection{Why Noetherian schemes preserve coherence under kernels}

Suppose
\[
f:\mathcal{F}\longrightarrow\mathcal{G}
\]
is a morphism of coherent sheaves on a Noetherian scheme \(X\).

Locally, on
\[
U=\operatorname{Spec} A,
\]
this becomes a homomorphism
\[
\phi:M\longrightarrow N
\]
of finitely generated \(A\)-modules.

Because \(A\) is Noetherian:
\[
\ker\phi
\]
is finitely generated,
\[
\operatorname{im}\phi
\]
is finitely generated, and
\[
\operatorname{coker}\phi
\]
is finitely generated.

Therefore:
\[
\boxed{
	\ker f,\qquad
	\operatorname{im} f,\qquad
	\operatorname{coker} f
	\text{ are coherent.}
}
\]

Without Noetherianity, a submodule of a finitely generated module need not be finitely generated. Consequently, kernels and images of morphisms between finite objects may fail to remain finite.

% ------------------------------------------------------------
\subsection{Pullback and pushforward of sheaves}

\subsubsection{Pullback as extension of scalars}

Let
\[
g:X\longrightarrow Y
\]
be a morphism of schemes.

For a sheaf \(\mathcal{F}\) of \(\mathcal{O}_Y\)-modules, the pullback is defined by
\[
g^*\mathcal{F}
=
g^{-1}\mathcal{F}
\otimes_{g^{-1}\mathcal{O}_Y}
\mathcal{O}_X.
\]

On affine charts, this construction is simply tensor product.

Suppose
\[
V=\operatorname{Spec} A\subseteq Y,
\qquad
U=\operatorname{Spec} B\subseteq g^{-1}(V),
\]
and the restricted morphism corresponds to a ring map
\[
A\longrightarrow B.
\]

If
\[
\mathcal{F}|_V\cong\widetilde{M},
\]
then
\[
(g^*\mathcal{F})|_U
\cong
\widetilde{M\otimes_A B}.
\]

Thus pullback is geometrically the same operation as extending scalars:
\[
M
\quad\rightsquigarrow\quad
M\otimes_A B.
\]

Since tensor products of modules are modules, pullback preserves quasi-coherence:
\[
\boxed{
	\mathcal{F}\text{ quasi-coherent}
	\quad\Longrightarrow\quad
	g^*\mathcal{F}\text{ quasi-coherent}.
}
\]

If \(M\) has a finite presentation
\[
A^p\longrightarrow A^q\longrightarrow M\longrightarrow 0,
\]
then tensoring with \(B\) gives
\[
B^p\longrightarrow B^q\longrightarrow M\otimes_A B\longrightarrow 0.
\]

Hence finite presentation is preserved under pullback. In the standard Noetherian setting of coherent sheaves, this implies:
\[
\boxed{
	\mathcal{F}\text{ coherent}
	\quad\Longrightarrow\quad
	g^*\mathcal{F}\text{ coherent}.
}
\]

\begin{remark}
	For arbitrary non-Noetherian schemes, coherence must be handled with care, because finite presentation and coherence need not coincide without additional hypotheses. In the usual locally Noetherian setting, the pullback statement above is exactly the required form.
\end{remark}

\subsubsection{Pushforward and the loss of finiteness}

The direct image or pushforward of a sheaf \(\mathcal{G}\) on \(X\) is defined by
\[
(g_*\mathcal{G})(V)
=
\mathcal{G}(g^{-1}(V))
\]
for open subsets
\[
V\subseteq Y.
\]

Unlike pullback, pushforward does not behave locally as a simple finite module operation. It collects all sections over inverse images of open sets, and these spaces of sections may be very large.

For example, consider
\[
p:\mathbf{A}^1_k\longrightarrow\operatorname{Spec} k.
\]

The coherent sheaf
\[
\mathcal{O}_{\mathbf{A}^1_k}
\]
has pushforward corresponding to
\[
\Gamma(\mathbf{A}^1_k,\mathcal{O}_{\mathbf{A}^1_k})
=
k[t].
\]

As a \(k\)-vector space,
\[
k[t]
\]
is infinite-dimensional. Hence it is not a finite \(k\)-module, and the corresponding sheaf on \(\operatorname{Spec} k\) is not coherent.

Thus:
\[
\boxed{
	\text{pushforward does not preserve coherence in general.}
}
\]

The underlying reason is that coherence expresses local finite generation, while global sections over a nonproper space may form an infinite module.

This also explains why properness is so important in algebraic geometry. A fundamental theorem states that for a proper morphism of Noetherian schemes, direct images of coherent sheaves are coherent. The failure in the affine-line example occurs because
\[
\mathbf{A}^1_k\longrightarrow\operatorname{Spec} k
\]
is not proper.

\subsubsection{Quasi-coherence of pushforwards}

Pushforward is better behaved with respect to quasi-coherence than with respect to coherence. Under standard hypotheses, for example when \(g\) is quasi-compact and quasi-separated, one has
\[
\mathcal{G}\text{ quasi-coherent}
\quad\Longrightarrow\quad
g_*\mathcal{G}\text{ quasi-coherent}.
\]

In the affine case this is immediate. If
\[
g:\operatorname{Spec} B\longrightarrow\operatorname{Spec} A
\]
corresponds to a ring map
\[
A\longrightarrow B
\]
and
\[
\mathcal{G}=\widetilde{N}
\]
for a \(B\)-module \(N\), then
\[
g_*\mathcal{G}
\]
is the quasi-coherent sheaf associated to the same abelian group \(N\), now viewed as an \(A\)-module by restriction of scalars.

However, even when this pushforward remains quasi-coherent, it need not be coherent, because \(N\) may fail to be finitely generated as an \(A\)-module.

% ------------------------------------------------------------
\subsection{Closed immersions and ideal sheaves}

\subsubsection{Closed subsets versus closed subschemes}

A closed subset of an affine scheme
\[
X=\operatorname{Spec} A
\]
is determined by a radical ideal:
\[
V(I)=V(\sqrt{I}).
\]

Thus the topological closed subset cannot distinguish between \(I\) and \(\sqrt I\).

A closed \emph{subscheme}, however, retains the entire ideal \(I\), including nilpotent information.

For example, inside
\[
\mathbf{A}^1_k=\operatorname{Spec} k[x],
\]
the ideals
\[
(x)
\qquad\text{and}\qquad
(x^2)
\]
have the same vanishing set:
\[
V(x)=V(x^2)=\{0\}.
\]

But they define different closed subschemes:
\[
\operatorname{Spec} k[x]/(x)
\qquad\text{and}\qquad
\operatorname{Spec} k[x]/(x^2).
\]

The first is an ordinary reduced point. The second is a nonreduced point with an infinitesimal thickening. Indeed, in
\[
k[x]/(x^2),
\]
the class of \(x\) is nonzero but satisfies
\[
x^2=0.
\]

Therefore closed subschemes encode more geometry than closed subsets.

\subsubsection{Definition of a closed immersion}

A morphism
\[
i:Z\longrightarrow X
\]
is a \emph{closed immersion} if:

\begin{enumerate}[label=\arabic*.]
	\item the map of topological spaces identifies \(Z\) homeomorphically with a closed subset of \(X\);
	\item the induced morphism of sheaves of rings
	\[
	i^\#:\mathcal{O}_X\longrightarrow i_*\mathcal{O}_Z
	\]
	is surjective.
\end{enumerate}

The kernel of this surjection is denoted
\[
\mathcal{I}_{Z/X}
=
\ker\!\left(
\mathcal{O}_X\longrightarrow i_*\mathcal{O}_Z
\right).
\]

There is therefore an exact sequence
\[
0
\longrightarrow
\mathcal{I}_{Z/X}
\longrightarrow
\mathcal{O}_X
\longrightarrow
i_*\mathcal{O}_Z
\longrightarrow
0.
\]

The sheaf \(\mathcal{I}_{Z/X}\) is called the \emph{ideal sheaf} defining \(Z\) inside \(X\).

\subsubsection{Why the kernel is a sheaf of ideals}

For every open subset
\[
U\subseteq X,
\]
one has
\[
\mathcal{I}_{Z/X}(U)
=
\ker\!\left(
\mathcal{O}_X(U)
\longrightarrow
\mathcal{O}_Z(U\cap Z)
\right).
\]

Since this is the kernel of a ring homomorphism,
\[
\mathcal{I}_{Z/X}(U)
\]
is an ideal of
\[
\mathcal{O}_X(U).
\]

Thus the ideal-sheaf property is simply the sheaf-theoretic version of the elementary algebraic fact that the kernel of a quotient map
\[
A\longrightarrow A/I
\]
is the ideal \(I\).

\subsubsection{The affine local form of a closed immersion}

Every closed immersion is affine locally of the form
\[
\operatorname{Spec}(A/I)\hookrightarrow\operatorname{Spec} A.
\]

Indeed, if
\[
U=\operatorname{Spec} A\subseteq X
\]
is affine and \(Z\hookrightarrow X\) is a closed immersion, then
\[
Z\cap U
\]
is affine and there exists an ideal
\[
I\subseteq A
\]
such that
\[
Z\cap U\cong\operatorname{Spec}(A/I).
\]

On this affine chart, the ideal sheaf is
\[
\mathcal{I}_{Z/X}|_U\cong\widetilde I.
\]

This immediately implies that
\[
\mathcal{I}_{Z/X}
\]
is quasi-coherent.

If \(X\) is Noetherian, then every ideal \(I\subseteq A\) is finitely generated. Hence
\[
\widetilde I
\]
is coherent on every affine open set, and therefore
\[
\mathcal{I}_{Z/X}
\]
is coherent.

Thus:
\[
\boxed{
	Z\hookrightarrow X\text{ closed immersion}
	\quad\Longrightarrow\quad
	\mathcal{I}_{Z/X}\text{ quasi-coherent},
}
\]
and, if \(X\) is Noetherian,
\[
\boxed{
	\mathcal{I}_{Z/X}\text{ coherent}.
}
\]

% ------------------------------------------------------------
\subsection{The correspondence between ideals and closed subschemes}

\subsubsection{The affine correspondence}

Let
\[
X=\operatorname{Spec} A.
\]

Given an ideal
\[
I\subseteq A,
\]
the quotient map
\[
A\longrightarrow A/I
\]
induces a closed immersion
\[
\operatorname{Spec}(A/I)\hookrightarrow\operatorname{Spec} A.
\]

Conversely, every closed immersion into \(\operatorname{Spec} A\) is isomorphic to one obtained in this way.

Thus there is a bijection:
\[
\boxed{
	\left\{
	\text{closed subschemes of }\operatorname{Spec} A
	\right\}
	\longleftrightarrow
	\left\{
	\text{ideals }I\subseteq A
	\right\}.
}
\]

The underlying closed subset of
\[
\operatorname{Spec}(A/I)
\]
is
\[
V(I)=V(\sqrt I).
\]

Thus the reduced closed subscheme associated to \(I\) is defined by \(\sqrt I\), while the possible nilpotent thickening is contained in the difference between \(I\) and \(\sqrt I\).

\subsubsection{The global correspondence}

On a general scheme \(X\), one replaces ideals of a single ring by quasi-coherent ideal sheaves
\[
\mathcal{I}\subseteq\mathcal{O}_X.
\]

Given a closed subscheme
\[
i:Z\hookrightarrow X,
\]
one obtains its ideal sheaf
\[
\mathcal{I}_{Z/X}
=
\ker(\mathcal{O}_X\to i_*\mathcal{O}_Z).
\]

Conversely, suppose
\[
\mathcal{I}\subseteq\mathcal{O}_X
\]
is a quasi-coherent ideal sheaf.

On an affine open subset
\[
U=\operatorname{Spec} A\subseteq X,
\]
there exists an ideal
\[
I_U\subseteq A
\]
such that
\[
\mathcal{I}|_U\cong\widetilde{I_U}.
\]

Define locally
\[
Z\cap U=\operatorname{Spec}(A/I_U).
\]

Because \(\mathcal{I}\) is a sheaf and its local ideals restrict compatibly, these affine closed subschemes agree on overlaps and glue to a global closed subscheme
\[
Z\hookrightarrow X.
\]

Thus:
\[
\boxed{
	\left\{
	\text{closed subschemes of }X
	\right\}
	\longleftrightarrow
	\left\{
	\text{quasi-coherent ideal sheaves }
	\mathcal{I}\subseteq\mathcal{O}_X
	\right\}.
}
\]

\subsubsection{The structure sheaf of a closed subscheme}

If a closed subscheme \(Z\hookrightarrow X\) is defined by an ideal sheaf
\[
\mathcal{I}\subseteq\mathcal{O}_X,
\]
then its structure sheaf is locally the quotient
\[
\mathcal{O}_X/\mathcal{I}.
\]

More precisely, one has
\[
i_*\mathcal{O}_Z\cong \mathcal{O}_X/\mathcal{I}.
\]

Hence the defining exact sequence is
\[
\boxed{
	0
	\longrightarrow
	\mathcal{I}
	\longrightarrow
	\mathcal{O}_X
	\longrightarrow
	i_*\mathcal{O}_Z
	\longrightarrow
	0.
}
\]

This exact sequence is one of the basic constructions in scheme theory. It allows geometric questions about the subscheme \(Z\) to be studied by relating its structure sheaf to the ambient scheme \(X\).

For example, infinitesimal information about how \(Z\) sits inside \(X\) is encoded by the quotient
\[
\mathcal{I}/\mathcal{I}^2,
\]
called the \emph{conormal sheaf} of \(Z\) in \(X\).

Its dual, when locally free, is the normal bundle:
\[
\mathcal{N}_{Z/X}
=
\mathcal{H}om_{\mathcal{O}_Z}
\left(
\mathcal{I}/\mathcal{I}^2,
\mathcal{O}_Z
\right).
\]

Thus ideal sheaves are not only devices for defining closed subschemes; they are also the starting point for tangent, normal, and infinitesimal geometry.

% ------------------------------------------------------------
\subsection{Conceptual comparison of the main ideas}

The four groups of results studied above fit together as follows.

\begin{center}
	\renewcommand{\arraystretch}{1.25}
	\begin{tabular}{
			>{\raggedright\arraybackslash}p{3.7cm}
			>{\raggedright\arraybackslash}p{5.1cm}
			>{\raggedright\arraybackslash}p{5.1cm}
		}
		\toprule
		\textbf{Topic}
		&
		\textbf{Local algebraic object}
		&
		\textbf{Global geometric meaning}
		\\
		\midrule
		
		Divisor class groups
		&
		Height-one primes and valuations in normal domains
		&
		Failure of codimension-one subvarieties to be globally principal
		\\[4pt]
		
		Morphisms \(\mathbf{P}^n\to\mathbf{P}^m\)
		&
		Base-point-free homogeneous polynomials of one degree
		&
		Rigidity of projective space and preservation of full dimension
		\\[4pt]
		
		Quasi-coherent sheaves
		&
		\(A\)-modules on \(\operatorname{Spec}(A)\)
		&
		Sheaves that can be studied by affine algebra
		\\[4pt]
		
		Coherent sheaves
		&
		Finitely generated or finitely presented modules
		&
		Finite geometric data stable under suitable operations
		\\[4pt]
		
		Closed subschemes
		&
		Ideals \(I\subseteq A\) and quotients \(A/I\)
		&
		Closed sets together with nilpotent and infinitesimal structure
		\\
		\bottomrule
	\end{tabular}
\end{center}

The repeated pattern is:

\[
\boxed{
	\text{geometric construction on schemes}
	\quad\longleftrightarrow\quad
	\text{algebraic construction on affine coordinate rings}.
}
\]

In the divisor-class problems, removing divisors reduces a global calculation to units in a factorial ring.

In the projective-morphism problem, pulling back the hyperplane bundle reduces a morphism to homogeneous polynomials and their common zero locus.

In the sheaf problems, kernels, images, cokernels, and pullbacks reduce locally to the corresponding operations on modules.

In the closed-subscheme problem, a geometric embedding is exactly a quotient construction
\[
A\longrightarrow A/I
\]
performed compatibly on affine charts.

These are basic examples of the general philosophy of scheme theory: schemes are global geometric spaces built from rings, and the most effective proofs move systematically between these two perspectives.

% ============================================================
\section{Examples for Problems 4--7:
	Divisor Classes, Projective Morphisms, Coherent Sheaves, and Closed Subschemes}

The following examples illustrate the four main themes developed in the preceding theory section:

\[
\operatorname{Cl}(X),\qquad
\mathbf P^n\longrightarrow \mathbf P^m,\qquad
\text{quasi-coherent and coherent sheaves},\qquad
\text{closed subschemes and ideal sheaves}.
\]

The examples begin with elementary affine schemes and then proceed to the projective surfaces, affine cones, morphisms, and ideal sheaves occurring directly in the problems.

% ------------------------------------------------------------
\subsection{Examples of divisor class groups}

\subsubsection{\texorpdfstring{The affine line: \(\operatorname{Cl}(\mathbf A^1_k)=0\)}{The affine line}}

\textbf{Example.}

Let
\[
X=\mathbf A^1_k=\operatorname{Spec}k[t].
\]

Since \(k[t]\) is a principal ideal domain, it is in particular a unique factorization domain. Every height-one prime ideal is generated by one irreducible polynomial. Hence every prime divisor is principal, and therefore
\[
\boxed{
	\operatorname{Cl}(\mathbf A^1_k)=0.
}
\]

Assume for simplicity that \(k\) is algebraically closed. Then the closed points are
\[
V(t-a),
\qquad a\in k.
\]

Each point is principal:
\[
\operatorname{div}(t-a)=[V(t-a)].
\]

For example, the divisor
\[
2[V(t)]-3[V(t-1)]
\]
is principal because
\[
\operatorname{div}\left(\frac{t^2}{(t-1)^3}\right)
=
2[V(t)]-3[V(t-1)].
\]

Thus every formal combination of closed points on \(\mathbf A^1_k\) is represented by the zeros and poles of a rational function.

\begin{remark}
	If \(k\) is not algebraically closed, the codimension-one points correspond to irreducible polynomials in \(k[t]\), rather than only to linear polynomials \(t-a\). The conclusion
	\[
	\operatorname{Cl}(\mathbf A^1_k)=0
	\]
	remains unchanged because \(k[t]\) is still a UFD.
\end{remark}

% ------------------------------------------------------------
\subsubsection{\texorpdfstring{The affine plane: \(\operatorname{Cl}(\mathbf A^2_k)=0\)}{The affine plane}}

\textbf{Example.}

Let
\[
X=\mathbf A^2_k=\operatorname{Spec}k[x,y].
\]

Prime divisors on \(X\) correspond to irreducible polynomial equations
\[
f(x,y)=0.
\]

Examples include
\[
V(x),\qquad
V(y),\qquad
V(x-y),\qquad
V(x^2+y^2-1),
\]
provided the defining polynomial is irreducible over \(k\).

Because
\[
k[x,y]
\]
is a unique factorization domain, every irreducible codimension-one closed subset is principal. Hence
\[
\boxed{
	\operatorname{Cl}(\mathbf A^2_k)=0.
}
\]

For instance, the divisor
\[
D=2V(x)-V(y)-V(x-y)
\]
is principal:
\[
D
=
\operatorname{div}\left(\frac{x^2}{y(x-y)}\right).
\]

This example explains why affine spaces are useful open subsets in divisor-class computations: once a scheme is reduced to an affine factorial open part, all remaining class-group information must lie on the removed boundary.

% ------------------------------------------------------------
\subsubsection{\texorpdfstring{The projective line: \(\operatorname{Cl}(\mathbf P^1_k)\cong \mathbf Z\)}{The projective line}}

\textbf{Example.}

Let
\[
X=\mathbf P^1_k.
\]

Choose the point at infinity
\[
\infty=[1:0].
\]

Its complement is
\[
\mathbf P^1_k\setminus\{\infty\}\cong\mathbf A^1_k.
\]

Since
\[
\operatorname{Cl}(\mathbf A^1_k)=0,
\]
the class of the point at infinity generates the class group:
\[
\mathbf Z[\infty]\twoheadrightarrow\operatorname{Cl}(\mathbf P^1_k).
\]

On the affine chart with coordinate
\[
t=\frac{X_0}{X_1},
\]
the rational function \(t-a\), for \(a\in k\), has a zero at the point \(a\) and a pole at infinity:
\[
\operatorname{div}(t-a)
=
[a]-[\infty].
\]

Therefore every \(k\)-rational point is linearly equivalent to \(\infty\):
\[
[a]\sim[\infty].
\]

More generally, the degree of a divisor
\[
D=\sum_i n_i[p_i]
\]
is
\[
\deg D=\sum_i n_i[k(p_i):k].
\]

Every principal divisor has degree zero. On \(\mathbf P^1_k\), degree classifies divisor classes, giving
\[
\boxed{
	\operatorname{Cl}(\mathbf P^1_k)\cong\mathbf Z.
}
\]

For example,
\[
[0]+[1]-2[\infty]
\]
is principal because
\[
\operatorname{div}\bigl(t(t-1)\bigr)
=
[0]+[1]-2[\infty].
\]

Thus, unlike the affine line, the projective line has a nontrivial class group: the missing point at infinity supplies one free divisor class.

% ------------------------------------------------------------
\subsubsection{\texorpdfstring{The surface \(\mathbf P^1_k\times\mathbf P^1_k\)}{The surface P1 times P1}}

\textbf{Example.}

Let
\[
X=\mathbf P^1_k\times\mathbf P^1_k.
\]

There are two natural types of irreducible curves:
\[
F_1=\{p\}\times\mathbf P^1_k,
\qquad
F_2=\mathbf P^1_k\times\{q\}.
\]

These are the two families of rulings of the surface.

Choose
\[
D_1=\{\infty\}\times\mathbf P^1_k,
\qquad
D_2=\mathbf P^1_k\times\{\infty\}.
\]

Removing them gives
\[
X\setminus(D_1\cup D_2)
\cong
\mathbf A^1_k\times\mathbf A^1_k
\cong
\mathbf A^2_k.
\]

Since
\[
\operatorname{Cl}(\mathbf A^2_k)=0,
\]
the classes \([D_1]\) and \([D_2]\) generate \(\operatorname{Cl}(X)\).

Suppose there were a relation
\[
a[D_1]+b[D_2]=0.
\]

Then
\[
aD_1+bD_2=\operatorname{div}(f)
\]
for some rational function \(f\). Since the divisor of \(f\) is supported entirely on \(D_1\cup D_2\), the restriction of \(f\) to
\[
\mathbf A^2_k
\]
must be a unit. But
\[
k[s,t]^\times=k^\times.
\]

Hence \(f\) must be constant, so
\[
\operatorname{div}(f)=0.
\]

Therefore
\[
a=b=0.
\]

Consequently,
\[
\boxed{
	\operatorname{Cl}(\mathbf P^1_k\times\mathbf P^1_k)
	\cong
	\mathbf Z\oplus\mathbf Z.
}
\]

A divisor class may be represented by a pair
\[
(a,b).
\]

The class
\[
(1,0)
\]
is represented by a divisor of the form
\[
\{p\}\times\mathbf P^1_k,
\]
while the class
\[
(0,1)
\]
is represented by a divisor of the form
\[
\mathbf P^1_k\times\{q\}.
\]

For example, the diagonal
\[
\Delta
=
\left\{
([S_0:S_1],[T_0:T_1])
:
S_0T_1-S_1T_0=0
\right\}
\]
is defined by an equation of bidegree \((1,1)\). Therefore
\[
[\Delta]=(1,1).
\]

Likewise, a bihomogeneous equation of bidegree \((2,3)\) defines a divisor whose class is
\[
(2,3).
\]

% ------------------------------------------------------------
\subsubsection{Principal divisors on \(\mathbf P^1_k\times\mathbf P^1_k\)}

\textbf{Example.}

Choose affine coordinates
\[
s=\frac{S_0}{S_1},
\qquad
t=\frac{T_0}{T_1}
\]
on the dense affine open subset
\[
\mathbf A^1_k\times\mathbf A^1_k
\subseteq
\mathbf P^1_k\times\mathbf P^1_k.
\]

Regarded as a rational function on the complete surface, \(s\) vanishes along
\[
\{0\}\times\mathbf P^1_k
\]
and has a pole along
\[
\{\infty\}\times\mathbf P^1_k.
\]

Hence
\[
\operatorname{div}(s)
=
(\{0\}\times\mathbf P^1_k)
-
(\{\infty\}\times\mathbf P^1_k).
\]

Thus all vertical rulings are linearly equivalent.

Similarly,
\[
\operatorname{div}(t)
=
(\mathbf P^1_k\times\{0\})
-
(\mathbf P^1_k\times\{\infty\}),
\]
so all horizontal rulings are linearly equivalent.

However, no rational function gives a relation identifying a vertical ruling with a horizontal ruling. Thus the two directions remain independent in the class group.

% ------------------------------------------------------------
\subsubsection{\texorpdfstring{The quadric cone \(xy=z^2\)}{The quadric cone xy=z2}}

\textbf{Example.}

Let
\[
X=\operatorname{Spec}A,
\qquad
A=k[x,y,z]/(xy-z^2).
\]

Consider the prime divisor
\[
D=V(x,z).
\]

Its coordinate ring is
\[
A/(x,z)
\cong
k[y],
\]
so
\[
D\cong\mathbf A^1_k.
\]

Now consider the open subset
\[
U=D(x)\subseteq X.
\]

Since \(x\) is invertible on \(U\), the defining equation gives
\[
y=\frac{z^2}{x}.
\]

Therefore
\[
\Gamma(U,\mathcal O_U)
\cong
k[x,x^{-1},z].
\]

This is a factorial ring, so
\[
\operatorname{Cl}(U)=0.
\]

The complement \(V(x)\) has only one irreducible codimension-one component, namely \(D\). Thus \([D]\) generates \(\operatorname{Cl}(X)\).

At the generic point of \(D\), the element \(y\) is invertible, and
\[
x=\frac{z^2}{y}.
\]

Hence \(x\) vanishes twice along \(D\):
\[
\operatorname{ord}_D(x)=2.
\]

Therefore
\[
\operatorname{div}(x)=2D.
\]

Consequently,
\[
2[D]=0.
\]

The divisor \(D\) itself is not principal, so its class is nonzero. Therefore
\[
\boxed{
	\operatorname{Cl}(X)\cong\mathbf Z/2\mathbf Z.
}
\]

There is another ruling line
\[
D'=V(y,z).
\]

Since
\[
\operatorname{div}(z)=D+D',
\]
we obtain
\[
[D']=-[D].
\]

But in a group of order two,
\[
-[D]=[D].
\]

Therefore both ruling lines represent the same nontrivial divisor class:
\[
[D']=[D].
\]

% ------------------------------------------------------------
\subsubsection{\texorpdfstring{The family \(xy=z^n\)}{The family xy=zn}}

\textbf{Example.}

The preceding example extends naturally to the family
\[
X_n
=
\operatorname{Spec}
k[x,y,z]/(xy-z^n),
\qquad n\geq 2.
\]

Let
\[
D=V(x,z).
\]

On the open subset
\[
D(x)\subseteq X_n,
\]
one has
\[
y=\frac{z^n}{x},
\]
and therefore
\[
\Gamma(D(x),\mathcal O_{X_n})
\cong
k[x,x^{-1},z].
\]

Thus the complement is factorial, and \([D]\) generates the class group.

At the generic point of \(D\), the element \(y\) is a unit, and
\[
x=\frac{z^n}{y}.
\]

Therefore
\[
\operatorname{div}(x)=nD.
\]

Hence
\[
n[D]=0.
\]

The same unit calculation as before shows that this is the only relation. Consequently,
\[
\boxed{
	\operatorname{Cl}(X_n)\cong\mathbf Z/n\mathbf Z.
}
\]

In particular,
\[
xy=z^2
\quad\Longrightarrow\quad
\operatorname{Cl}(X_2)\cong\mathbf Z/2\mathbf Z,
\]
\[
xy=z^3
\quad\Longrightarrow\quad
\operatorname{Cl}(X_3)\cong\mathbf Z/3\mathbf Z,
\]
and
\[
xy=z^4
\quad\Longrightarrow\quad
\operatorname{Cl}(X_4)\cong\mathbf Z/4\mathbf Z.
\]

This family gives a useful collection of normal affine surface singularities with finite cyclic divisor class groups.

% ------------------------------------------------------------
\subsubsection{\texorpdfstring{The three-dimensional quadric cone \(xy=zw\)}{The quadric cone xy=zw}}

\textbf{Example.}

Let
\[
X=\operatorname{Spec}A,
\qquad
A=k[x,y,z,w]/(xy-zw).
\]

On the open subset
\[
U=D(x),
\]
one may solve for \(y\):
\[
y=\frac{zw}{x}.
\]

Hence
\[
\Gamma(U,\mathcal O_U)
\cong
k[x,x^{-1},z,w].
\]

This is factorial, so
\[
\operatorname{Cl}(U)=0.
\]

The complement \(V(x)\) has two irreducible codimension-one components:
\[
D=V(x,z),
\qquad
E=V(x,w).
\]

Indeed,
\[
A/(x)
\cong
k[y,z,w]/(zw),
\]
whose minimal primes are \((z)\) and \((w)\).

Their coordinate rings are
\[
A/(x,z)\cong k[y,w],
\]
and
\[
A/(x,w)\cong k[y,z].
\]

Therefore
\[
D\cong\mathbf A^2_k,
\qquad
E\cong\mathbf A^2_k.
\]

At the generic point of \(D\), the functions \(y\) and \(w\) are invertible, and
\[
x=\frac{w}{y}z.
\]

Hence
\[
\operatorname{ord}_D(x)=1.
\]

Similarly, at the generic point of \(E\),
\[
x=\frac{z}{y}w,
\]
so
\[
\operatorname{ord}_E(x)=1.
\]

Therefore
\[
\operatorname{div}(x)=D+E.
\]

Consequently,
\[
[D]+[E]=0.
\]

This is the only relation between the two generators, and hence
\[
\boxed{
	\operatorname{Cl}(X)\cong\mathbf Z.
}
\]

One may take \([D]\) as a generator, with
\[
[E]=-[D].
\]

% ------------------------------------------------------------
\subsubsection{Comparison of the class-group examples}

The examples above may be summarized as follows:

\[
\begin{array}{|p{3.3cm}|p{4.0cm}|p{3.5cm}|p{3.0cm}|}
	\hline
	\textbf{Scheme}
	&
	\textbf{Factorial open subset}
	&
	\textbf{Principal relation}
	&
	\textbf{Class group}
	\\
	\hline
	\(\mathbf A^1_k\)
	&
	\(\mathbf A^1_k\)
	&
	\(\text{Every divisor is principal}\)
	&
	\(0\)
	\\
	\hline
	\(\mathbf A^2_k\)
	&
	\(\mathbf A^2_k\)
	&
	\(\text{Every divisor is principal}\)
	&
	\(0\)
	\\
	\hline
	\(\mathbf P^1_k\)
	&
	\(\mathbf A^1_k\)
	&
	\([a]-[\infty]=\operatorname{div}(t-a)\)
	&
	\(\mathbf Z\)
	\\
	\hline
	\(\mathbf P^1_k\times\mathbf P^1_k\)
	&
	\(\mathbf A^2_k\)
	&
	\(\text{No relation between the two rulings}\)
	&
	\(\mathbf Z^2\)
	\\
	\hline
	\(\operatorname{Spec}k[x,y,z]/(xy-z^2)\)
	&
	\(\operatorname{Spec}k[x,x^{-1},z]\)
	&
	\(\operatorname{div}(x)=2D\)
	&
	\(\mathbf Z/2\mathbf Z\)
	\\
	\hline
	\(\operatorname{Spec}k[x,y,z]/(xy-z^n)\)
	&
	\(\operatorname{Spec}k[x,x^{-1},z]\)
	&
	\(\operatorname{div}(x)=nD\)
	&
	\(\mathbf Z/n\mathbf Z\)
	\\
	\hline
	\(\operatorname{Spec}k[x,y,z,w]/(xy-zw)\)
	&
	\(\operatorname{Spec}k[x,x^{-1},z,w]\)
	&
	\(\operatorname{div}(x)=D+E\)
	&
	\(\mathbf Z\)
	\\
	\hline
\end{array}
\]

The contrast between the two cone examples is especially important:
\[
\operatorname{div}(x)=2D
\]
produces a torsion class, while
\[
\operatorname{div}(x)=D+E
\]
produces one relation between two free generators.

% ------------------------------------------------------------
\subsection{Examples of morphisms between projective spaces}

\subsubsection{A constant morphism}

\textbf{Example.}

Define
\[
\varphi:\mathbf P^2_k\longrightarrow\mathbf P^3_k
\]
by
\[
\varphi([X:Y:Z])=[1:0:0:0].
\]

Every point is sent to the same image, so
\[
\varphi(\mathbf P^2_k)=\{[1:0:0:0]\}.
\]

Thus
\[
\dim\varphi(\mathbf P^2_k)=0.
\]

The pullback of the hyperplane line bundle is trivial:
\[
\varphi^*\mathcal O_{\mathbf P^3}(1)
\cong
\mathcal O_{\mathbf P^2}.
\]

This is the degree-zero case in the general theory of morphisms from projective space.

% ------------------------------------------------------------
\subsubsection{A linear embedding \(\mathbf P^2_k\hookrightarrow\mathbf P^3_k\)}

\textbf{Example.}

Define
\[
\varphi:\mathbf P^2_k\longrightarrow\mathbf P^3_k
\]
by
\[
\varphi([X:Y:Z])=[X:Y:Z:0].
\]

The image is the projective plane
\[
V(W)\subseteq\mathbf P^3_k.
\]

Therefore
\[
\varphi(\mathbf P^2_k)\cong\mathbf P^2_k,
\]
and
\[
\dim\varphi(\mathbf P^2_k)=2.
\]

The morphism is given by linear forms, and therefore
\[
\varphi^*\mathcal O_{\mathbf P^3}(1)
\cong
\mathcal O_{\mathbf P^2}(1).
\]

This is a basic example of a nonconstant morphism whose image has the full dimension of the source.

% ------------------------------------------------------------
\subsubsection{\texorpdfstring{A degree-\(d\) morphism \(\mathbf P^1_k\to\mathbf P^1_k\)}{A degree d morphism P1 to P1}}

\textbf{Example.}

For an integer \(d\geq1\), define
\[
\varphi_d:\mathbf P^1_k\longrightarrow\mathbf P^1_k
\]
by
\[
\varphi_d([S:T])
=
[S^d:T^d].
\]

The two homogeneous polynomials \(S^d\) and \(T^d\) never vanish simultaneously on \(\mathbf P^1_k\). Hence this is a well-defined morphism.

On the affine chart \(T\neq0\), with coordinate
\[
u=\frac ST,
\]
the map is
\[
u\longmapsto u^d.
\]

The pullback of the hyperplane bundle is
\[
\varphi_d^*\mathcal O_{\mathbf P^1}(1)
\cong
\mathcal O_{\mathbf P^1}(d).
\]

The image is all of \(\mathbf P^1_k\), so
\[
\dim\varphi_d(\mathbf P^1_k)=1.
\]

For \(d=2\), the map
\[
[S:T]\longmapsto[S^2:T^2]
\]
is a degree-two morphism. It may identify different points, but it does not lower the dimension of the image.

% ------------------------------------------------------------
\subsubsection{The Veronese embedding of \(\mathbf P^1_k\) as a conic}

\textbf{Example.}

Define
\[
\nu_2:\mathbf P^1_k\longrightarrow\mathbf P^2_k
\]
by
\[
\nu_2([S:T])
=
[S^2:ST:T^2].
\]

The three homogeneous coordinates never vanish simultaneously, so this defines a morphism everywhere on \(\mathbf P^1_k\).

Write the homogeneous coordinates on \(\mathbf P^2_k\) as
\[
[X:Y:Z].
\]

For every point in the image,
\[
XZ-Y^2
=
(S^2)(T^2)-(ST)^2
=
0.
\]

Thus
\[
\nu_2(\mathbf P^1_k)
\subseteq
V(XZ-Y^2).
\]

Conversely, the smooth conic
\[
C=V(XZ-Y^2)\subseteq\mathbf P^2_k
\]
is parametrized by this map. Therefore
\[
\nu_2(\mathbf P^1_k)=C.
\]

The image is one-dimensional:
\[
\dim C=1.
\]

Moreover,
\[
\nu_2^*\mathcal O_{\mathbf P^2}(1)
\cong
\mathcal O_{\mathbf P^1}(2).
\]

This example shows that a nonconstant projective morphism may transform a projective line into a curved-looking embedded conic while preserving its dimension.

% ------------------------------------------------------------
\subsubsection{The Veronese embedding of \(\mathbf P^2_k\)}

\textbf{Example.}

Define
\[
\nu_2:\mathbf P^2_k\longrightarrow\mathbf P^5_k
\]
by
\[
\nu_2([X:Y:Z])
=
[X^2:XY:XZ:Y^2:YZ:Z^2].
\]

The six coordinates are homogeneous of degree \(2\), and they have no common zero in \(\mathbf P^2_k\). Hence \(\nu_2\) is a morphism.

Its image is a projective surface in \(\mathbf P^5_k\), called the Veronese surface. Since \(\nu_2\) is an embedding,
\[
\dim\nu_2(\mathbf P^2_k)=2.
\]

The associated pullback relation is
\[
\nu_2^*\mathcal O_{\mathbf P^5}(1)
\cong
\mathcal O_{\mathbf P^2}(2).
\]

Again, the image of a nonconstant morphism from projective space has the full dimension of the source.

% ------------------------------------------------------------
\subsubsection{\texorpdfstring{A finite morphism \(\mathbf P^2_k\to\mathbf P^2_k\)}{A finite morphism P2 to P2}}

\textbf{Example.}

For \(d\geq1\), define
\[
\varphi_d:\mathbf P^2_k\longrightarrow\mathbf P^2_k
\]
by
\[
\varphi_d([X:Y:Z])
=
[X^d:Y^d:Z^d].
\]

The homogeneous coordinates \(X^d,Y^d,Z^d\) have no common projective zero, so this is a morphism.

Its image is all of \(\mathbf P^2_k\), and therefore
\[
\dim\varphi_d(\mathbf P^2_k)=2.
\]

For \(d>1\), this morphism generally fails to be injective. Nevertheless, its image still has full dimension.

Thus the conclusion of the projective-morphism theorem is not that every nonconstant map is an embedding. Rather, it says that a nonconstant morphism from \(\mathbf P^n_k\) cannot have a positive-dimensional image of dimension strictly less than \(n\).

% ------------------------------------------------------------
\subsubsection{\texorpdfstring{Why \([X:Y]\) does not define a morphism \(\mathbf P^2_k\to\mathbf P^1_k\)}{An invalid projection}}

\textbf{Example.}

One might try to define a map
\[
\psi:\mathbf P^2_k\dashrightarrow\mathbf P^1_k
\]
by
\[
\psi([X:Y:Z])=[X:Y].
\]

This expression is valid whenever
\[
(X,Y)\neq(0,0).
\]

However, at the point
\[
[0:0:1]\in\mathbf P^2_k,
\]
both proposed target coordinates vanish:
\[
[X:Y]=[0:0].
\]

But \([0:0]\) is not a point of projective space. Therefore the formula is not defined at \([0:0:1]\).

Thus
\[
[X:Y:Z]\longmapsto[X:Y]
\]
defines only a rational map
\[
\mathbf P^2_k\dashrightarrow\mathbf P^1_k,
\]
not a morphism on all of \(\mathbf P^2_k\).

This gives a concrete explanation of the general theorem:
\[
\boxed{
	\text{there is no nonconstant everywhere-defined morphism }
	\mathbf P^2_k\longrightarrow\mathbf P^1_k.
}
\]

The point \([0:0:1]\) is called a base point of the proposed linear system.

% ------------------------------------------------------------
\subsubsection{Higher-degree coordinates do not automatically remove base points}

\textbf{Example.}

Consider the formula
\[
[X:Y:Z]
\longmapsto
[X^2:XY:Y^2].
\]

All three target coordinates have the same degree. However, at the point
\[
[0:0:1],
\]
they all vanish:
\[
X^2=XY=Y^2=0.
\]

Therefore this formula again gives only a rational map
\[
\mathbf P^2_k\dashrightarrow\mathbf P^2_k,
\]
not a morphism defined everywhere.

This illustrates the essential requirement for a projective morphism:

\[
\boxed{
	\text{homogeneous coordinates of one degree must have no common zero.}
}
\]

% ------------------------------------------------------------
\subsection{Examples of quasi-coherent and coherent sheaves}

\subsubsection{The structure sheaf of an affine scheme}

\textbf{Example.}

Let
\[
X=\operatorname{Spec}A.
\]

The structure sheaf is the sheaf associated to the \(A\)-module \(A\):
\[
\mathcal O_X\cong\widetilde A.
\]

Therefore \(\mathcal O_X\) is quasi-coherent.

If \(A\) is Noetherian, then \(A\) is finitely generated as a module over itself. Hence
\[
\mathcal O_X
\]
is coherent.

For example,
\[
\mathcal O_{\mathbf A^2_k}
\]
is coherent because
\[
\mathbf A^2_k=\operatorname{Spec}k[x,y]
\]
and \(k[x,y]\) is Noetherian.

% ------------------------------------------------------------
\subsubsection{The ideal sheaf of the origin in \(\mathbf A^2_k\)}

\textbf{Example.}

Let
\[
X=\mathbf A^2_k=\operatorname{Spec}k[x,y].
\]

The origin is defined by the maximal ideal
\[
I=(x,y).
\]

The associated ideal sheaf is
\[
\widetilde{(x,y)}\subseteq\mathcal O_{\mathbf A^2_k}.
\]

The quotient is
\[
k[x,y]/(x,y)\cong k,
\]
which is the coordinate ring of the reduced origin.

Thus there is an exact sequence
\[
0
\longrightarrow
\widetilde{(x,y)}
\longrightarrow
\mathcal O_{\mathbf A^2_k}
\longrightarrow
\mathcal O_{\{(0,0)\}}
\longrightarrow
0.
\]

Since the ideal \((x,y)\) is finitely generated, its associated sheaf is coherent.

This example shows concretely how a closed point is described by a coherent ideal sheaf.

% ------------------------------------------------------------
\subsubsection{Multiplication by \(t\) on the affine line}

\textbf{Example.}

Let
\[
X=\mathbf A^1_k=\operatorname{Spec}k[t].
\]

Consider the morphism
\[
\mathcal O_X
\xrightarrow{\cdot t}
\mathcal O_X.
\]

On modules, this is the map
\[
k[t]\xrightarrow{\cdot t}k[t].
\]

Since \(k[t]\) is a domain,
\[
\ker(\cdot t)=0.
\]

Its image is the ideal
\[
(t)\subseteq k[t].
\]

Its cokernel is
\[
k[t]/(t)\cong k.
\]

Therefore, as sheaves,
\[
\ker(\cdot t)=0,
\]
\[
\operatorname{im}(\cdot t)=\widetilde{(t)},
\]
and
\[
\operatorname{coker}(\cdot t)
\cong
\mathcal O_{\{0\}}.
\]

Hence one obtains the exact sequence
\[
0
\longrightarrow
\mathcal O_{\mathbf A^1_k}
\xrightarrow{\cdot t}
\mathcal O_{\mathbf A^1_k}
\longrightarrow
\mathcal O_{\{0\}}
\longrightarrow
0.
\]

This sequence says that when regular functions are taken modulo those divisible by \(t\), only their value at the origin remains.

% ------------------------------------------------------------
\subsubsection{The ideal of the origin as an image}

\textbf{Example.}

Let
\[
X=\mathbf A^2_k=\operatorname{Spec}k[x,y].
\]

Define a morphism
\[
\varphi:\mathcal O_X^{\oplus 2}\longrightarrow\mathcal O_X
\]
by
\[
\varphi(a,b)=ax+by.
\]

On global modules, this is
\[
k[x,y]^2\longrightarrow k[x,y],
\qquad
(a,b)\longmapsto ax+by.
\]

Its image is precisely the ideal
\[
(x,y).
\]

Thus
\[
\operatorname{im}\varphi=\widetilde{(x,y)}.
\]

The cokernel is
\[
k[x,y]/(x,y)\cong k,
\]
and consequently
\[
\operatorname{coker}\varphi
\cong
\mathcal O_{\{(0,0)\}}.
\]

The kernel is generated by the relation
\[
(-y,x),
\]
since
\[
(-y)x+xy=0.
\]

Therefore there is an exact sequence
\[
0
\longrightarrow
\mathcal O_{\mathbf A^2_k}
\xrightarrow{
	\begin{pmatrix}
		-y\\
		x
	\end{pmatrix}
}
\mathcal O_{\mathbf A^2_k}^{\oplus 2}
\xrightarrow{
	\begin{pmatrix}
		x & y
	\end{pmatrix}
}
\mathcal O_{\mathbf A^2_k}
\longrightarrow
\mathcal O_{\{(0,0)\}}
\longrightarrow
0.
\]

This is the Koszul resolution of the structure sheaf of the origin in \(\mathbf A^2_k\).

% ------------------------------------------------------------
\subsubsection{A quasi-coherent sheaf that is not coherent}

\textbf{Example.}

Let
\[
X=\operatorname{Spec}k.
\]

A quasi-coherent sheaf on \(X\) is simply a \(k\)-vector space, viewed as a sheaf on the one-point scheme.

Consider the \(k\)-vector space
\[
M=k[t].
\]

The associated sheaf
\[
\widetilde{k[t]}
\]
is quasi-coherent.

However,
\[
k[t]
\]
is infinite-dimensional as a vector space over \(k\). Hence it is not finitely generated as a \(k\)-module.

Therefore
\[
\boxed{
	\widetilde{k[t]}
	\text{ is quasi-coherent but not coherent on }\operatorname{Spec}k.
}
\]

This is exactly the sheaf that occurs as the pushforward
\[
p_*\mathcal O_{\mathbf A^1_k}
\]
under the structure morphism
\[
p:\mathbf A^1_k\longrightarrow\operatorname{Spec}k.
\]

% ------------------------------------------------------------
\subsubsection{A coherent skyscraper sheaf}

\textbf{Example.}

Let
\[
X=\mathbf A^1_k=\operatorname{Spec}k[t].
\]

The origin is defined by the ideal
\[
(t).
\]

Its structure sheaf, regarded as a sheaf on \(X\), is
\[
\mathcal O_{\{0\}}
\cong
\widetilde{k[t]/(t)}.
\]

Since
\[
k[t]/(t)\cong k
\]
is finitely generated as a \(k[t]\)-module, the sheaf
\[
\mathcal O_{\{0\}}
\]
is coherent.

Likewise,
\[
\widetilde{k[t]/(t^2)}
\]
is coherent, but it describes a nonreduced thickened point rather than the ordinary reduced point.

Thus coherence does not require a scheme to be reduced. It expresses finite generation, not the absence of nilpotents.

% ------------------------------------------------------------
\subsection{Examples of pullbacks and pushforwards}

\subsubsection{Pulling back the origin under the squaring map}

\textbf{Example.}

Let
\[
g:\mathbf A^1_k\longrightarrow\mathbf A^1_k
\]
be defined by
\[
u\longmapsto u^2.
\]

Writing the coordinate on the target as \(t\), this morphism corresponds to the ring homomorphism
\[
k[t]\longrightarrow k[u],
\qquad
t\longmapsto u^2.
\]

Let
\[
Z=V(t)\subseteq\mathbf A^1_k
\]
be the origin in the target. Its structure sheaf, as an \(\mathcal O_{\mathbf A^1_k}\)-module, is
\[
\mathcal O_Z\cong\widetilde{k[t]/(t)}.
\]

Its pullback is
\[
g^*\mathcal O_Z
\cong
\widetilde{
	\bigl(k[t]/(t)\bigr)\otimes_{k[t]}k[u]
}.
\]

Since \(t\) acts on \(k[u]\) as multiplication by \(u^2\), one obtains
\[
\bigl(k[t]/(t)\bigr)\otimes_{k[t]}k[u]
\cong
k[u]/(u^2).
\]

Therefore
\[
\boxed{
	g^*\mathcal O_Z
	\cong
	\widetilde{k[u]/(u^2)}.
}
\]

Set-theoretically, the inverse image of \(0\) is only the point \(u=0\). Scheme-theoretically, however, the inverse image is the double point
\[
\operatorname{Spec}k[u]/(u^2).
\]

Thus the pullback remembers the multiplicity of the zero of the function \(u^2\).

% ------------------------------------------------------------
\subsubsection{Pulling back the origin under a cubic map}

\textbf{Example.}

Let
\[
g:\mathbf A^1_k\longrightarrow\mathbf A^1_k
\]
be defined by
\[
t\longmapsto u^3.
\]

Equivalently, on coordinate rings,
\[
k[t]\longrightarrow k[u],
\qquad
t\longmapsto u^3.
\]

Pulling back the origin \(V(t)\) gives
\[
k[u]/(u^3).
\]

Thus the scheme-theoretic inverse image is
\[
\operatorname{Spec}k[u]/(u^3),
\]
a triple point at the origin.

The underlying topological inverse image is still just
\[
\{0\},
\]
but its scheme structure records multiplicity three.

% ------------------------------------------------------------
\subsubsection{Pullback of \(\mathcal O_{\mathbf P^1}(1)\) under a degree-\(d\) map}

\textbf{Example.}

Let
\[
\varphi_d:\mathbf P^1_k\longrightarrow\mathbf P^1_k
\]
be defined by
\[
\varphi_d([S:T])=[S^d:T^d].
\]

Then
\[
\varphi_d^*\mathcal O_{\mathbf P^1}(1)
\cong
\mathcal O_{\mathbf P^1}(d).
\]

For instance, if
\[
d=2,
\]
then
\[
\varphi_2^*\mathcal O_{\mathbf P^1}(1)
\cong
\mathcal O_{\mathbf P^1}(2).
\]

This may also be seen directly from divisors. The point
\[
[0:1]\in\mathbf P^1_k
\]
is defined by the equation
\[
S=0.
\]

Under the degree-two map, its defining equation pulls back to
\[
S^2=0.
\]

Thus the pullback divisor has multiplicity two:
\[
\varphi_2^*([0:1])=2[0:1].
\]

More generally,
\[
\varphi_d^*([0:1])=d[0:1].
\]

% ------------------------------------------------------------
\subsubsection{A pushforward that is not coherent}

\textbf{Example.}

Consider the structure morphism
\[
p:\mathbf A^1_k\longrightarrow\operatorname{Spec}k.
\]

The sheaf
\[
\mathcal O_{\mathbf A^1_k}
\]
is coherent.

However,
\[
p_*\mathcal O_{\mathbf A^1_k}
\]
corresponds to the \(k\)-vector space
\[
\Gamma(\mathbf A^1_k,\mathcal O_{\mathbf A^1_k})
=
k[t].
\]

As a vector space over \(k\),
\[
k[t]
=
k\oplus kt\oplus kt^2\oplus kt^3\oplus\cdots
\]
is infinite-dimensional.

A coherent sheaf on \(\operatorname{Spec}k\) corresponds to a finite-dimensional \(k\)-vector space. Therefore
\[
\boxed{
	p_*\mathcal O_{\mathbf A^1_k}
	\text{ is not coherent.}
}
\]

It remains quasi-coherent, because it is still the sheaf associated to the \(k\)-module \(k[t]\).

% ------------------------------------------------------------
\subsubsection{A positive pushforward example: a finite morphism}

\textbf{Example.}

Let
\[
g:\mathbf A^1_k\longrightarrow\mathbf A^1_k
\]
be defined by
\[
u\longmapsto u^2.
\]

Write the coordinate on the target as \(t\). The corresponding ring map is
\[
k[t]\longrightarrow k[u],
\qquad
t\longmapsto u^2.
\]

As a module over \(k[t]\), the ring \(k[u]\) is generated by
\[
1
\qquad\text{and}\qquad
u.
\]

Indeed, every polynomial \(F(u)\in k[u]\) can be written uniquely in the form
\[
F(u)=A(u^2)+uB(u^2)
\]
for some polynomials \(A,B\in k[t]\). Hence
\[
k[u]\cong k[t]\oplus u\,k[t]
\]
as a \(k[t]\)-module.

Thus \(k[u]\) is finite over \(k[t]\), and therefore
\[
g_*\mathcal O_{\mathbf A^1_k}
\]
is coherent.

This should be contrasted with
\[
\mathbf A^1_k\longrightarrow\operatorname{Spec}k,
\]
where the pushforward of the structure sheaf is not coherent.

% ------------------------------------------------------------
\subsection{Examples of closed subschemes and ideal sheaves}

\subsubsection{A reduced point in \(\mathbf A^1_k\)}

\textbf{Example.}

Let
\[
X=\mathbf A^1_k=\operatorname{Spec}k[x].
\]

The ideal
\[
I=(x)
\]
defines the closed subscheme
\[
Z=\operatorname{Spec}k[x]/(x).
\]

Since
\[
k[x]/(x)\cong k,
\]
we have
\[
Z\cong\operatorname{Spec}k.
\]

This is the ordinary reduced point at the origin.

Its ideal sheaf is
\[
\mathcal I_{Z/X}=\widetilde{(x)}.
\]

The defining exact sequence is
\[
0
\longrightarrow
\widetilde{(x)}
\longrightarrow
\mathcal O_{\mathbf A^1_k}
\longrightarrow
\mathcal O_{\{0\}}
\longrightarrow
0.
\]

% ------------------------------------------------------------
\subsubsection{A double point in \(\mathbf A^1_k\)}

\textbf{Example.}

Now take the ideal
\[
I=(x^2).
\]

It defines the closed subscheme
\[
Z=\operatorname{Spec}k[x]/(x^2).
\]

Its underlying topological space is still the single point \(0\), because
\[
V(x^2)=V(x).
\]

However, its coordinate ring contains a nonzero nilpotent element:
\[
x\neq0,
\qquad
x^2=0
\quad\text{in }k[x]/(x^2).
\]

Therefore \(Z\) is not reduced. It is a double point, or first infinitesimal thickening of the origin.

The reduced point and the double point have the same underlying topological subset but different scheme structures:
\[
\operatorname{Spec}k[x]/(x)
\not\cong
\operatorname{Spec}k[x]/(x^2).
\]

% ------------------------------------------------------------
\subsubsection{A triple point in \(\mathbf A^1_k\)}

\textbf{Example.}

The ideal
\[
I=(x^3)
\]
defines
\[
Z=\operatorname{Spec}k[x]/(x^3).
\]

As a \(k\)-vector space,
\[
k[x]/(x^3)
=
k\oplus kx\oplus kx^2.
\]

In this ring,
\[
x^3=0,
\qquad
x^2\neq0.
\]

Thus \(Z\) is a triple infinitesimal thickening of the origin.

There is a chain of closed subschemes:
\[
\operatorname{Spec}k[x]/(x)
\subseteq
\operatorname{Spec}k[x]/(x^2)
\subseteq
\operatorname{Spec}k[x]/(x^3)
\subseteq
\mathbf A^1_k.
\]

Each has the same underlying point, but each successive scheme remembers more nilpotent information.

% ------------------------------------------------------------
\subsubsection{Two reduced points in \(\mathbf A^1_k\)}

\textbf{Example.}

Let
\[
I=(x(x-1))\subseteq k[x].
\]

Then
\[
Z=\operatorname{Spec}k[x]/(x(x-1)).
\]

Since the ideals
\[
(x)
\qquad\text{and}\qquad
(x-1)
\]
are comaximal, the Chinese remainder theorem gives
\[
k[x]/(x(x-1))
\cong
k[x]/(x)\times k[x]/(x-1)
\cong
k\times k.
\]

Therefore \(Z\) is the disjoint union of two reduced points:
\[
Z=\{0,1\}.
\]

This is very different from
\[
\operatorname{Spec}k[x]/(x^2),
\]
which has only one topological point but possesses nilpotent structure.

% ------------------------------------------------------------
\subsubsection{A line in \(\mathbf A^2_k\)}

\textbf{Example.}

Let
\[
X=\mathbf A^2_k=\operatorname{Spec}k[x,y].
\]

The ideal
\[
I=(y)
\]
defines the \(x\)-axis:
\[
Z=V(y).
\]

Its coordinate ring is
\[
k[x,y]/(y)\cong k[x].
\]

Hence
\[
Z\cong\mathbf A^1_k.
\]

Its ideal sheaf is
\[
\mathcal I_{Z/X}=\widetilde{(y)}.
\]

The corresponding exact sequence is
\[
0
\longrightarrow
\widetilde{(y)}
\longrightarrow
\mathcal O_{\mathbf A^2_k}
\longrightarrow
\mathcal O_{V(y)}
\longrightarrow
0.
\]

% ------------------------------------------------------------
\subsubsection{A doubled line in \(\mathbf A^2_k\)}

\textbf{Example.}

Replace the ideal \((y)\) by
\[
I=(y^2).
\]

The corresponding closed subscheme is
\[
Z=V(y^2).
\]

Its underlying topological set is still the \(x\)-axis:
\[
V(y^2)=V(y).
\]

However, its coordinate ring is
\[
k[x,y]/(y^2),
\]
in which the class of \(y\) is nonzero but nilpotent:
\[
y\neq0,
\qquad
y^2=0.
\]

Thus \(Z\) is a doubled line. It contains first-order infinitesimal information in the direction normal to the ordinary line \(V(y)\).

% ------------------------------------------------------------
\subsubsection{The union of the coordinate axes}

\textbf{Example.}

In
\[
\mathbf A^2_k=\operatorname{Spec}k[x,y],
\]
consider the ideal
\[
I=(xy).
\]

The corresponding closed subscheme is
\[
Z=V(xy)=V(x)\cup V(y).
\]

Thus \(Z\) is the union of the two coordinate axes.

Its coordinate ring is
\[
k[x,y]/(xy).
\]

The scheme is reduced, because the ideal \((xy)\) is radical in the polynomial ring over a field:
\[
(xy)=(x)\cap(y).
\]

However, it is reducible, since it is the union of the two irreducible components
\[
V(x)
\qquad\text{and}\qquad
V(y).
\]

It is also singular at the origin, where the two components meet.

This example distinguishes three notions:

\[
\boxed{
	\text{reduced does not imply irreducible, and irreducible does not imply smooth.}
}
\]

% ------------------------------------------------------------
\subsubsection{A fat point in \(\mathbf A^2_k\)}

\textbf{Example.}

Let
\[
X=\mathbf A^2_k=\operatorname{Spec}k[x,y].
\]

The ordinary origin is defined by
\[
(x,y).
\]

Now consider the square of this ideal:
\[
I=(x,y)^2=(x^2,xy,y^2).
\]

The resulting closed subscheme is
\[
Z=\operatorname{Spec}
k[x,y]/(x^2,xy,y^2).
\]

As a \(k\)-vector space, its coordinate ring has basis
\[
1,\quad x,\quad y.
\]

The classes of \(x\) and \(y\) are nonzero, but every second-order expression vanishes:
\[
x^2=xy=y^2=0.
\]

Thus \(Z\) is a fat point supported at the origin. It remembers first-order infinitesimal directions in both coordinate directions.

The reduced origin
\[
\operatorname{Spec}k[x,y]/(x,y)
\]
remembers only the point itself, whereas
\[
\operatorname{Spec}k[x,y]/(x,y)^2
\]
also remembers the first-order tangent information around that point.

% ------------------------------------------------------------
\subsubsection{The diagonal in \(\mathbf A^1_k\times\mathbf A^1_k\)}

\textbf{Example.}

Let
\[
X=\mathbf A^1_k\times\mathbf A^1_k
=
\operatorname{Spec}k[x,y].
\]

The diagonal consists of the points satisfying
\[
x=y.
\]

It is therefore defined by the ideal
\[
I=(x-y).
\]

Its coordinate ring is
\[
k[x,y]/(x-y)\cong k[x].
\]

Hence the diagonal is isomorphic to the affine line:
\[
\Delta\cong\mathbf A^1_k.
\]

Its ideal sheaf is
\[
\mathcal I_{\Delta/X}
=
\widetilde{(x-y)}.
\]

The diagonal is a fundamental example in scheme theory. For a general scheme \(X\), the diagonal morphism
\[
\Delta_X:X\longrightarrow X\times X
\]
plays a central role in the definition of separated morphisms and in the construction of differentials.

% ------------------------------------------------------------
\subsubsection{The graph of a polynomial map}

\textbf{Example.}

Consider the morphism
\[
f:\mathbf A^1_k\longrightarrow\mathbf A^1_k,
\qquad
x\longmapsto x^2.
\]

Its graph is the subset of
\[
\mathbf A^1_k\times\mathbf A^1_k
=
\operatorname{Spec}k[x,y]
\]
defined by
\[
y=x^2.
\]

Thus the graph is the closed subscheme
\[
\Gamma_f=V(y-x^2),
\]
defined by the ideal
\[
I=(y-x^2).
\]

Its coordinate ring is
\[
k[x,y]/(y-x^2)\cong k[x].
\]

Therefore
\[
\Gamma_f\cong\mathbf A^1_k.
\]

This example shows that morphisms can often be studied geometrically by studying their graphs as closed subschemes of products.

% ------------------------------------------------------------
\subsubsection{A hyperplane in projective space}

\textbf{Example.}

Let
\[
H=V(X_0)\subseteq\mathbf P^n_k.
\]

Then
\[
H\cong\mathbf P^{n-1}_k.
\]

The ideal sheaf of \(H\) is
\[
\mathcal I_H\cong\mathcal O_{\mathbf P^n}(-1).
\]

Indeed, multiplication by the defining homogeneous linear form \(X_0\) gives the exact sequence
\[
0
\longrightarrow
\mathcal O_{\mathbf P^n}(-1)
\xrightarrow{\cdot X_0}
\mathcal O_{\mathbf P^n}
\longrightarrow
\mathcal O_H
\longrightarrow
0.
\]

More generally, if a hypersurface
\[
Y=V(F)\subseteq\mathbf P^n_k
\]
is defined by a homogeneous polynomial \(F\) of degree \(d\), then
\[
\mathcal I_Y\cong\mathcal O_{\mathbf P^n}(-d),
\]
and there is an exact sequence
\[
0
\longrightarrow
\mathcal O_{\mathbf P^n}(-d)
\xrightarrow{\cdot F}
\mathcal O_{\mathbf P^n}
\longrightarrow
\mathcal O_Y
\longrightarrow
0.
\]

For example, the plane conic
\[
C=V(XZ-Y^2)\subseteq\mathbf P^2_k
\]
has defining equation of degree \(2\), so
\[
\mathcal I_C\cong\mathcal O_{\mathbf P^2}(-2),
\]
and
\[
0
\longrightarrow
\mathcal O_{\mathbf P^2}(-2)
\xrightarrow{\cdot(XZ-Y^2)}
\mathcal O_{\mathbf P^2}
\longrightarrow
\mathcal O_C
\longrightarrow
0.
\]

% ------------------------------------------------------------
\subsection{Examples connecting the different topics}

\subsubsection{\texorpdfstring{The divisor \(D=V(x,z)\) on the cone \(xy=z^2\)}{A divisor on the quadric cone}}

\textbf{Example.}

Let
\[
X=\operatorname{Spec}k[x,y,z]/(xy-z^2).
\]

The closed subscheme
\[
D=V(x,z)
\]
is defined by the ideal
\[
I_D=(x,z).
\]

Its coordinate ring is
\[
\frac{k[x,y,z]/(xy-z^2)}{(x,z)}
\cong
k[y].
\]

Hence
\[
D\cong\mathbf A^1_k.
\]

Thus \(D\) is simultaneously:

\begin{itemize}
	\item a closed subscheme of \(X\), defined by the ideal sheaf associated to \((x,z)\);
	\item an irreducible codimension-one subvariety of \(X\), hence a prime divisor;
	\item a representative of the nontrivial class in
	\[
	\operatorname{Cl}(X)\cong\mathbf Z/2\mathbf Z.
	\]
\end{itemize}

The relation
\[
\operatorname{div}(x)=2D
\]
says that \(D\) is not itself principal, but twice \(D\) is principal.

This example joins together closed subschemes, prime divisors, rational functions, and divisor classes.

% ------------------------------------------------------------
\subsubsection{\texorpdfstring{The two divisors on \(xy=zw\)}{Two divisors on the three-dimensional cone}}

\textbf{Example.}

Let
\[
X=\operatorname{Spec}k[x,y,z,w]/(xy-zw).
\]

Consider the two closed subschemes
\[
D=V(x,z),
\qquad
E=V(x,w).
\]

They are defined by the ideals
\[
I_D=(x,z),
\qquad
I_E=(x,w).
\]

Their coordinate rings are
\[
\frac{k[x,y,z,w]/(xy-zw)}{(x,z)}
\cong
k[y,w],
\]
and
\[
\frac{k[x,y,z,w]/(xy-zw)}{(x,w)}
\cong
k[y,z].
\]

Therefore
\[
D\cong\mathbf A^2_k,
\qquad
E\cong\mathbf A^2_k.
\]

Moreover, the zero locus of \(x\) decomposes as
\[
V(x)=D\cup E.
\]

At the divisor level,
\[
\operatorname{div}(x)=D+E.
\]

Therefore, in the divisor class group,
\[
[D]+[E]=0.
\]

Hence
\[
[E]=-[D].
\]

This gives a direct connection between:

\begin{itemize}
	\item decomposition of a closed subscheme into irreducible components;
	\item principal divisors of regular functions;
	\item relations inside the divisor class group.
\end{itemize}

% ------------------------------------------------------------
\subsubsection{The conic as both an image and a closed subscheme}

\textbf{Example.}

Consider the Veronese morphism
\[
\nu_2:\mathbf P^1_k\longrightarrow\mathbf P^2_k,
\qquad
[S:T]\longmapsto[S^2:ST:T^2].
\]

Its image is the conic
\[
C=V(XZ-Y^2)\subseteq\mathbf P^2_k.
\]

Thus \(C\) can be described in two complementary ways:

\begin{enumerate}[label=\arabic*.]
	\item As the image of a morphism from projective space:
	\[
	C=\nu_2(\mathbf P^1_k).
	\]
	
	\item As the closed subscheme of \(\mathbf P^2_k\) defined by the homogeneous ideal
	\[
	(XZ-Y^2).
	\]
\end{enumerate}

Its ideal sheaf satisfies
\[
\mathcal I_C\cong\mathcal O_{\mathbf P^2}(-2),
\]
and the defining exact sequence is
\[
0
\longrightarrow
\mathcal O_{\mathbf P^2}(-2)
\xrightarrow{\cdot(XZ-Y^2)}
\mathcal O_{\mathbf P^2}
\longrightarrow
\mathcal O_C
\longrightarrow
0.
\]

At the same time, the projective morphism satisfies
\[
\nu_2^*\mathcal O_{\mathbf P^2}(1)
\cong
\mathcal O_{\mathbf P^1}(2).
\]

Thus this single example simultaneously illustrates:

\begin{itemize}
	\item morphisms from projective space;
	\item homogeneous equations defining closed subschemes;
	\item coherent ideal sheaves;
	\item pullbacks of line bundles.
\end{itemize}

% ------------------------------------------------------------
\subsubsection{Scheme-theoretic inverse images and multiplicity}

\textbf{Example.}

Let
\[
g:\mathbf A^1_k\longrightarrow\mathbf A^1_k,
\qquad
u\longmapsto u^2.
\]

Let
\[
Z=V(t)
\]
be the origin in the target affine line.

Set-theoretically,
\[
g^{-1}(\{0\})=\{0\}.
\]

However, scheme-theoretically, the defining ideal pulls back by
\[
(t)\longmapsto(u^2).
\]

Therefore
\[
g^{-1}(Z)
=
\operatorname{Spec}k[u]/(u^2).
\]

Thus the scheme-theoretic inverse image is a double point.

Similarly, for the map
\[
u\longmapsto u^3,
\]
the inverse image of the origin is
\[
\operatorname{Spec}k[u]/(u^3),
\]
a triple point.

Hence:
\[
\boxed{
	\text{scheme-theoretic inverse images remember multiplicity.}
}
\]

This is one of the most important reasons that schemes retain ideal-theoretic information rather than only underlying point sets.

% ------------------------------------------------------------
\subsection{Summary table of the examples}

\renewcommand{\arraystretch}{1.35}
\begin{center}
	\begin{tabular}{|p{4.0cm}|p{5.2cm}|p{5.0cm}|}
		\hline
		\textbf{Example}
		&
		\textbf{Construction}
		&
		\textbf{Main conclusion}
		\\
		\hline
		\(\mathbf A^1_k\)
		&
		\(\operatorname{Spec}k[t]\)
		&
		\(\operatorname{Cl}(\mathbf A^1_k)=0\)
		\\
		\hline
		\(\mathbf A^2_k\)
		&
		\(\operatorname{Spec}k[x,y]\)
		&
		\(\operatorname{Cl}(\mathbf A^2_k)=0\)
		\\
		\hline
		\(\mathbf P^1_k\)
		&
		Remove the point \(\infty\)
		&
		\(\operatorname{Cl}(\mathbf P^1_k)\cong\mathbf Z\)
		\\
		\hline
		\(\mathbf P^1_k\times\mathbf P^1_k\)
		&
		Remove two ruling divisors
		&
		\(\operatorname{Cl}\cong\mathbf Z^2\)
		\\
		\hline
		\(xy=z^2\)
		&
		\(\operatorname{div}(x)=2D\)
		&
		\(\operatorname{Cl}\cong\mathbf Z/2\mathbf Z\)
		\\
		\hline
		\(xy=z^n\)
		&
		\(\operatorname{div}(x)=nD\)
		&
		\(\operatorname{Cl}\cong\mathbf Z/n\mathbf Z\)
		\\
		\hline
		\(xy=zw\)
		&
		\(\operatorname{div}(x)=D+E\)
		&
		\(\operatorname{Cl}\cong\mathbf Z\)
		\\
		\hline
		Constant projective map
		&
		\([X:Y:Z]\mapsto[1:0:0:0]\)
		&
		Image is a point
		\\
		\hline
		Linear embedding
		&
		\([X:Y:Z]\mapsto[X:Y:Z:0]\)
		&
		Image has dimension \(2\)
		\\
		\hline
		Veronese conic
		&
		\([S:T]\mapsto[S^2:ST:T^2]\)
		&
		Image is \(V(XZ-Y^2)\)
		\\
		\hline
		Invalid projection
		&
		\([X:Y:Z]\dashrightarrow[X:Y]\)
		&
		Not defined at \([0:0:1]\)
		\\
		\hline
		Origin in \(\mathbf A^2_k\)
		&
		Ideal \((x,y)\)
		&
		Reduced closed point
		\\
		\hline
		Double point
		&
		Ideal \((x^2)\subseteq k[x]\)
		&
		Nonreduced thickening of one point
		\\
		\hline
		Fat point
		&
		Ideal \((x,y)^2\subseteq k[x,y]\)
		&
		First-order infinitesimal directions at the origin
		\\
		\hline
		Doubled line
		&
		Ideal \((y^2)\subseteq k[x,y]\)
		&
		Nonreduced thickening of a line
		\\
		\hline
		Diagonal
		&
		Ideal \((x-y)\subseteq k[x,y]\)
		&
		Closed copy of \(\mathbf A^1_k\)
		\\
		\hline
		Pullback under \(u\mapsto u^2\)
		&
		\((t)\mapsto(u^2)\)
		&
		Reduced point pulls back to a double point
		\\
		\hline
		Noncoherent pushforward
		&
		\(\mathbf A^1_k\to\operatorname{Spec}k\)
		&
		\(p_*\mathcal O_{\mathbf A^1_k}\) is not coherent
		\\
		\hline
		Finite pushforward
		&
		\(u\mapsto u^2\) on \(\mathbf A^1_k\)
		&
		\(g_*\mathcal O_{\mathbf A^1_k}\) is coherent
		\\
		\hline
	\end{tabular}
\end{center}

% ------------------------------------------------------------
\subsection{Concluding observations}

The divisor examples demonstrate that the divisor class group detects subtle failures of global principality.

For the surface cone
\[
xy=z^2,
\]
the relation
\[
\operatorname{div}(x)=2D
\]
creates a nontrivial torsion class:
\[
[D]\neq0,
\qquad
2[D]=0.
\]

For the three-dimensional cone
\[
xy=zw,
\]
the relation
\[
\operatorname{div}(x)=D+E
\]
identifies two prime-divisor classes as opposites:
\[
[E]=-[D],
\]
leaving one free generator.

The projective examples show that a genuine morphism
\[
\mathbf P^n_k\longrightarrow\mathbf P^m_k
\]
must be represented by homogeneous polynomials with no common zero. The failed projection
\[
[X:Y:Z]\dashrightarrow[X:Y]
\]
is not a morphism precisely because of its base point. By contrast, the Veronese formula
\[
[S:T]\longmapsto[S^2:ST:T^2]
\]
has no base point and gives a genuine closed embedding.

The sheaf examples show concretely that operations on quasi-coherent sheaves are module-theoretic on affine schemes. The exact sequence
\[
0
\longrightarrow
\mathcal O_{\mathbf A^1_k}
\xrightarrow{\cdot t}
\mathcal O_{\mathbf A^1_k}
\longrightarrow
\mathcal O_{\{0\}}
\longrightarrow
0
\]
is simply the sheaf version of
\[
0
\longrightarrow
k[t]
\xrightarrow{\cdot t}
k[t]
\longrightarrow
k[t]/(t)
\longrightarrow
0.
\]

Finally, the closed-subscheme examples explain why scheme theory is richer than ordinary point-set geometry. Although
\[
V(x)=V(x^2)
\]
as subsets of \(\mathbf A^1_k\), the corresponding schemes
\[
\operatorname{Spec}k[x]/(x)
\qquad\text{and}\qquad
\operatorname{Spec}k[x]/(x^2)
\]
are not the same: the second contains a nilpotent infinitesimal direction.

Thus the examples support the central scheme-theoretic principle:
\[
\boxed{
	\text{geometric spaces are understood through their local rings, modules, ideals, and rational functions.}
}
\]

% ============================================================
% Add the following packages and TikZ libraries to the preamble
% if they are not already present:
%
% \usepackage{tikz}
% \usetikzlibrary{
	%     arrows.meta,
	%     positioning,
	%     calc,
	%     fit,
	%     shapes.geometric,
	%     shapes.multipart,
	%     decorations.pathreplacing,
	%     matrix
	% }
% \usepackage{booktabs}
% \usepackage{tabularx}
% \usepackage{array}
% \usepackage{adjustbox}
% \usepackage{xcolor}
%
% The following section uses the macros already introduced earlier:
% \mathbf{A}, \mathbf{P}, \mathbb{Z}, \operatorname{Cl}, \operatorname{Spec}, \mathcal{O}, \operatorname{div}.
% ============================================================

